\chapter{Safety without identification: transient amplification and a maximal pre-trade guard}\label{ch:p2}

    This chapter, the second paper of the thesis, condenses the paper \emph{From Human Markets to Machine Markets: Transient Amplification, an Exact Depth Condition, and a Maximal Pre-Trade Guard}, joint with Suboh Alkhushayni, Yuhao Chen, Mansi Bhavsar and Prashant Niroula. It keeps the theorems the rest of the thesis uses, with their proofs, and the empirical program in summary, with every registered test and its outcome in Table~\ref{p2:tab:alltests}. The full version, with every table, robustness check and the source of every number, is in the replication package (\texttt{thesis/ch3\_paper2\_full.tex}).

    \begin{quote}\small\textbf{Summary.} Eigenvalues and volatility do not bound the losses a market produces before it settles. In a limit order book with adaptive traders, cumulative transient risk is $S=\tilde\kappa/(m-q\tilde\kappa)$, long-run impact over reversion net of momentum; a book overshoots its direct impact if and only if $\varphi\varepsilon>(1-a)\kappa$; and a pre-trade guard on the summed impulse response is the largest safe admission rule and does not depend on how the traders split into cohorts. The summed response is certified with $\Theta(\sigma^2\epsilon^{-2}\log(1/\epsilon))$ experiments whatever the cohort spacing, while identifying the cohorts costs polynomially more, and the certified guard is minimax-optimal in lost capacity up to a logarithm. Of 25 registered empirical tests, 15 pass: six-hour estimates of the statistic forecast the size of the next hour's adverse excursion in thin Binance perpetuals beyond a full volatility baseline and VPIN, more weakly in thin US stocks, and not on Coinbase or in the deepest book; tests of a liquidation mechanism fail.\end{quote}

    \section{Introduction}

    On October 10, 2025, crypto derivatives markets liquidated an estimated \$19 billion of leveraged positions in a single day \citep{p2:shares2025}; on May 6, 2010, US equity prices fell and recovered by several percent within minutes \citep{p2:kirilenko2017}. In both episodes the losses that mattered were transient: positions were closed out on the path, before prices settled. Algorithmic markets are closed-loop systems, and a common shortcut equates their stability with eigenvalues inside the unit circle. For a leveraged intermediary that is the wrong object: a market can be spectrally stable and still push a dealer through its maintenance margin within seconds.

    The chapter asks three questions. When does a limit order book amplify a flow shock? How can a dealer or a venue control transient loss without identifying who is trading, which Chapter~\ref{ch:p1} shows is expensive from public data? And what information certifies a safe trading decision, at what cost? The answers, in a linear book model, are these.

    \paragraph{Spectral stability does not bound transient loss.} The safe shock radius of a stable market is the loss budget over its peak transient gain, and that gain can be arbitrarily large at any fixed spectral radius (Theorem~\ref{p2:thm:radius}).

    \paragraph{When a book amplifies.} A book with depletion and refill overshoots its direct impact if and only if $\varphi\varepsilon>(1-a)\kappa$; momentum alone cannot make a stable book overshoot; and cumulative risk has the closed form $S=\tilde\kappa/(m-q\tilde\kappa)$, which needs no depletion (Theorems~\ref{p2:thm:peak} and~\ref{p2:thm:sum}).

    \paragraph{The maximal guard.} Under a declared risk contract, worst-case loss is inventory times the shock bound times the summed absolute impulse response, so a guard on that sum admits exactly the safe actions (Lemma~\ref{p2:thm:guard}); a guard on the peak response admits unsafe ones. The guard survives estimation error (Theorem~\ref{p2:thm:robust}) and needs no identification of the traders (Proposition~\ref{p2:prop:invariance}).

    \paragraph{Certifying is cheaper than identifying.} Two markets can agree on their early responses and differ threefold in risk (Example~\ref{p2:ex:witness}). A finite window certifies the summed response under a memory bound and certifies nothing without one (Theorem~\ref{p2:thm:window}); with noisy experiments the cost of certification is $\Theta(\sigma^2\epsilon^{-2}\log(1/\epsilon))$ whatever the cohort structure, against $(\Delta\eta)^{-(4K-2)}$ for identification (Theorem~\ref{p2:thm:certcost}). The guard built on the certificate has a closed form and is minimax-optimal (Theorem~\ref{p2:thm:robustguard}); disclosure is worth most when memory is long, and a noisy proxy can be worth less than nothing (Theorem~\ref{p2:thm:disclosure}).

    \paragraph{Who controls the risk.} Reversion is a public good: dealers under-provide it, momentum traders over-provide feedback, and the corrective rebate is a function of the guard's own statistic (Proposition~\ref{p2:prop:equil}).

    \paragraph{Evidence.} Twenty-five primaries were registered; twenty-four were tested and one was not opened under its fixed sequence; fifteen pass (Table~\ref{p2:tab:alltests}). The deepest book, BTCUSDT, sits on the stable side of the amplification condition and fails a registered forecasting test. In thin books the statistic forecasts the size of the next hour's adverse excursion in every registered forecasting test on Binance, on new symbols and on untouched days, by up to $2.8$ points of $R^2$ beyond a full volatility baseline; about a fifth as much in thin US stocks; negligibly on Bybit; not at all on Coinbase. We have not identified a deeper mechanism: the effect does not run through open-interest liquidations, and it is carried by slow reversion relative to impact, which is what the statistic measures. As a risk control the certified guard lowers tail loss against a static limit but not against a volatility cap at matched exposure.

    \paragraph{What we do not claim.} The information-safety theorems concern declared observation contracts (controlled reset experiments, exact windows, audited data), not finite-sample guarantees for observational order-book records. Cohorts are any traders whose demand is a linear filter of past prices; nothing attributes an effect to algorithms rather than people. The forecasting result is reduced-form, and its size depends on the venue.

    \paragraph{Related work.} Models of transient impact and resilience \citep{p2:obizhaeva2013,p2:bouchaud2004,p2:bouchaud2018} take the decay of impact as given; we derive when depletion, refill and momentum make a book overshoot and give cumulative impact in closed form. Liquidity spirals \citep{p2:brunnermeier2009,p2:shleifer1992} need funding constraints; our amplification does not, and our guard is the control that bounds it. \citet{p2:sornette2023} show that non-normal interactions create transient bubbles. The guard's worst-case identity is the $\ell_1$ gain of $\ell_1$-optimal control \citep{p2:vidyasagar1986,p2:dahleh1987,p2:dahleh1995}; what is new is its use as an admission rule, its invariance to the cohort coordinates and its identification from prices. The early-warning benchmark is VPIN \citep{p2:easley2012,p2:andersen2014}; forecast evaluation uses \citet{p2:clark2007} with the volatility baselines of \citet{p2:corsi2009} and \citet{p2:riskmetrics1996}, and we pre-register following \citet{p2:nosek2018} and \citet{p2:harvey2017}. Rule 15c3-5 requires pre-trade controls on credit and capital \citep{p2:sec2010}; the guard is a candidate with a provable maximality property.

    \section{The closed-loop market}\label{p2:sec:model}

    \subsection{Price formation}
Time is discrete with bars of length $\Delta$. Let $p_t$ be the price, $p^*_t=p^*_{t-1}+\xi_t$ the unobserved fundamental and $x_t=p_t-p^*_t$ the pricing error. Competitive dealers post aggregate depth $D>0$, so price impact is $\kappa=1/D$, and fundamental arbitrage removes a fraction $\kappa\varphi_f$ of the error each period. With algorithmic demand $u_t$ and liquidity noise, the pricing error follows
\begin{equation}\label{p2:eq:price}
x_{t+1}=ax_t+\kappa u_t+w_{t+1},\qquad a=1-\kappa\varphi_f\in[0,1).
\end{equation}
Algorithmic cohorts trade $u_t=\sum_iq_iy_{i,t}$ with exponential memory $y_{i,t+1}=(1-\eta_i)y_{i,t}+\eta_ix_{t+1}$, a reduced form that arises when an adaptive desk updates its belief by steady-state Kalman filtering or regularized online gradient descent and maximizes expected continuation return $\beta_i y_{i,t}u_{i,t}$ minus quadratic execution cost $\tfrac12\lambda_{{\rm tc},i}u_{i,t}^2$ and inventory penalty $\tfrac12\gamma_{{\rm inv},i}I_{i,t}^2$ ($I_{i,t}=(1-\delta_{\rm inv})I_{i,t-1}+u_{i,t}$, $|I_{i,t}|\le\bar I_i$), yielding $u_{i,t}^*=q_iy_{i,t}-\psi_{{\rm inv},i}I_{i,t-1}$ with $q_i=\beta_i/(\lambda_{{\rm tc},i}+\gamma_{{\rm inv},i})$ and $\psi_{{\rm inv},i}=\gamma_{{\rm inv},i}(1-\delta_{\rm inv})/(\lambda_{{\rm tc},i}+\gamma_{{\rm inv},i})$ (Chapter~\ref{ch:p1}; this single-agent filtering-and-control rule is distinct from the joint dealer--trader equilibrium of Section~\ref{p2:sec:equil}). The baseline equations set $\psi_{{\rm inv},i}=0$, the limit in which inventory holding penalties are small relative to execution costs ($\gamma_{{\rm inv},i}\ll\lambda_{{\rm tc},i}$) or inventory is reset at the bar horizon; Proposition~\ref{p2:prop:inventory_S} below relaxes this and gives the closed-form summed response when cohorts carry inventory across bars with $\psi_{{\rm inv},i}>0$. Stacking $z_t=(x_t,y_{1,t},\ldots,y_{K,t})^\top$ gives $z_{t+1}=Jz_t+bw_{t+1}$, where $b=(1,\eta_1,\ldots,\eta_K)^\top$ because cohorts update on the same period's price (a unit of flow enters through $\kappa b$), and
\begin{equation}\label{p2:eq:J}
J=\begin{pmatrix} a & \kappa q_1 & \cdots & \kappa q_K\\ \eta_1 a & 1-\eta_1+\eta_1\kappa q_1 & \cdots & \eta_1\kappa q_K\\ \vdots & & \ddots & \vdots\\ \eta_K a & \eta_K\kappa q_1 & \cdots & 1-\eta_K+\eta_K\kappa q_K\end{pmatrix}.
\end{equation}

    \subsection{Stability and what an observer can use}
\begin{proposition}[Stability with one cohort]\label{p2:prop:jury}
With one cohort $(q,\eta)$ and $a=1-\kappa\varphi_f$, the market is Schur stable if and only if
\begin{equation}\label{p2:eq:jury}
q<\varphi_f,\qquad \kappa<\frac{4-2\eta}{(2-\eta)\varphi_f-\eta q}\ \ \text{when}\ (2-\eta)\varphi_f>\eta q,\qquad |1-\kappa\varphi_f|(1-\eta)<1 .
\end{equation}
Higher impact $\kappa$ (thinner depth) lowers price persistence $a$; feedback is destabilizing only when cohort aggressiveness $q$ reaches the arbitrage strength $\varphi_f$.
\end{proposition}
Chapter~\ref{ch:p1} shows that prices identify only the merged rate--weight measure $M=\sum_iq_i\eta_i\delta_{1-\eta_i}$ of the cohorts. A pre-trade control must therefore depend on the market only through objects that are the same for every cohort representation of $M$.

\begin{proposition}[Input--output invariance]\label{p2:prop:invariance}
Two populations with the same merged measure generate the same transfer function $H(z)=\sum_kg_kz^{-k}=\kappa z/[z-a-\kappa G(z)]$, $G(z)=\sum_iq_i\eta_iz/(z-1+\eta_i)$, and the same impulse coefficients $g_k=\kappa\,c_{\rm price}^\top J^kb$ for every $k\ge0$. Under any invertible change of state coordinates $z\mapsto Tz$, the Euclidean gain $\|CT^{-1}(TJT^{-1})^k\|_2$ depends on $T$, whereas the price path
\[
\Delta p_{t+k}=\sum_{j\ge1}g_{k-1+j}u_{t-j}+\sum_{m=0}^{k-1}g_{k-1-m}u_{t+m}
\]
and its worst case over $\|u\|_\infty\le\bar u$ depend only on $\{g_k\}$ and are invariant to both $T$ and cohort splitting.
\end{proposition}

Numerically, splitting a cohort of 1{,}000 random economies into two with the same memory rate leaves every impulse coefficient unchanged up to rounding ($\max_k|g_k(\mathcal E^A)-g_k(\mathcal E^B)|\le2.5\times10^{-12}$ over 200 lags; \texttt{recompute\_oct20.py}). Proposition~\ref{p2:prop:invariance} is the reason we build the guard on impulse coefficients and avoid state norms in latent coordinates. When a state norm is needed, we fix units by observable calm-market innovation standard deviations, which is what makes the spot statistic of Section~\ref{p2:sec:evidence} unit-free.
\begin{proposition}[Summed response with carried cohort inventory]\label{p2:prop:inventory_S}
Suppose each cohort $i$ carries inventory $I_{i,t}=(1-\delta_{\rm inv})I_{i,t-1}+u_{i,t}$ with $\delta_{\rm inv}\in(0,1]$ and submits order flow $u_{i,t}=q_i y_{i,t}-\psi_i I_{i,t-1}$ with $q_i\ge0$ and $\psi_i\ge0$. Then:
\begin{enumerate}[label=(\roman*)]
\item The closed-loop transfer function from a flow impulse to price is $H_{\rm inv}(z)=\kappa z/\bigl[z-a-\kappa G_{\rm inv}(z)\bigr]$ with
\[
G_{\rm inv}(z)=\sum_{i=1}^K \frac{q_i\eta_i z(z-1+\delta_{\rm inv})}{(z-1+\eta_i)(z-1+\delta_{\rm inv}+\psi_i)},
\]
which is invariant to splitting any cohort that shares $(\eta_i,\psi_i)$. Whenever the augmented state matrix is Schur stable, the signed cumulative price response is
\begin{equation}\label{p2:eq:Sinv}
S(\psi)=\sum_{t=0}^\infty h_t(\psi)=H_{\rm inv}(1)=\frac{\kappa}{1-a-\kappa\Phi_{\rm eff}(\psi)},\qquad \Phi_{\rm eff}(\psi)\equiv\sum_{i=1}^K\frac{\delta_{\rm inv}}{\delta_{\rm inv}+\psi_i}\,q_i\le\Phi.
\end{equation}
\item If $0\le\psi_i\le 1-\delta_{\rm inv}$ and $\delta_{\rm inv}\ge\eta_i$ for each $i$ and $a+\kappa\Phi_{\rm eff}(\psi)<1$, every open-loop feedback tap $k_m(\psi)$ and hence every closed-loop impulse coefficient $h_t(\psi)$ is nonnegative, so the worst-case $\ell_1$ gain equals the signed sum, $\|h(\psi)\|_1=\sum_{t=0}^\infty |h_t(\psi)|=S(\psi)$. When $\psi_i$ is large enough that inventory payback reverses the sign of late taps, $\|h(\psi)\|_1>S(\psi)$.
\end{enumerate}
\end{proposition}

Inventory mean-reversion $\psi_i>0$ acts as an automatic stabilizer that attenuates each cohort's long-run feedback weight from $q_i$ to $\frac{\delta_{\rm inv}}{\delta_{\rm inv}+\psi_i}q_i$ (across 500 random stable economies \eqref{p2:eq:Sinv} matches the simulated state-space impulse sum to $1.78\times10^{-15}$, and in the monotone regime $\|h(\psi)\|_1=S(\psi)$ to $1.78\times10^{-15}$, reducing $S$ from $2.7778$ at $\psi=0$ to $2.5000$ at the calibrated worked example; in the sign-reversal regime $\|h(\psi)\|_1=2.0960>S(\psi)=1.7241$; \texttt{recompute\_p2\_step2\_theory.py}).

    The model is linear: it describes the transient response to flow shocks of ordinary size, which is what a pre-trade control must bound, not regime changes such as exchange outages. Every input to the guard (impulse coefficients, an operator estimate and its confidence radius, inventory and margin) is observable to a dealer or a venue.

    \section{Spectral stability is not solvency}\label{p2:sec:radius}
    
\begin{theorem}[Safe radius, Kreiss bounds and the failure of the spectral radius]\label{p2:thm:radius}
Let $z_t=J^tz_0$ with $\rho(J)<1$, loss budget $B>0$ and loss output $C$.
\begin{enumerate}[label=(\roman*)]
\item The largest initial-state radius with $\sup_t\|CJ^tz_0\|_2\le B$ for every $\|z_0\|_2\le r$ is
\begin{equation}\label{p2:eq:rexact}
r^*_{\rm exact}=\frac{B}{M_C},\qquad M_C\equiv\sup_{t\ge0}\|CJ^t\|_2 .
\end{equation}
\item With $C=I$ and the Kreiss constant $\mathcal K(J)=\sup_{|z|>1}(|z|-1)\|(zI-J)^{-1}\|_2$, the Kreiss matrix theorem gives $\mathcal K(J)\le\sup_t\|J^t\|_2\le e\,n\,\mathcal K(J)$ for $J\in\R^{n\times n}$, hence
\begin{equation}\label{p2:eq:kreiss}
\frac{B}{e\,n\,\mathcal K(J)}\;\le\;r^*_{\rm exact}\;\le\;\frac{B}{\mathcal K(J)} .
\end{equation}
\item For any $\rho\in(0,1)$ and any $B>0$ there is a family $J_\theta$ with $\rho(J_\theta)=\rho$ for all $\theta$ and an output $C=(1,0)$ such that $r^*_{\rm exact}(J_\theta)\to0$ as $\theta\to\infty$.
\end{enumerate}
\end{theorem}

    For $J_\theta=\begin{psmallmatrix}0.8&\theta\\0&0.4\end{psmallmatrix}$ the spectral radius is $0.8$ throughout, while the price-row gain rises from 1 to 120 as $\theta$ goes from 0 to 100 and the safe radius falls by the same factor. Chapter~\ref{ch:p1} shows more: markets that share their eigenvalues, one-step price variance and stationary price variance can differ arbitrarily in worst-case transient gain, so neither eigenvalue nor volatility monitors bound transient loss. In the cohort economy with nonnegative weights the price response is a nonnegative mixture of decaying exponentials and peaks at the direct impact (Chapter~\ref{ch:p1}, positive-geometry lemma): non-normality is not sufficient for scalar price overshoot, and necessity for Euclidean state-norm growth does not extend to arbitrary input-output responses. In this positive cohort model, momentum alone does not supply overshoot. Depth depletion does.

    \section{When does a book amplify? An exact condition}\label{p2:sec:t6}

    Transient path risk takes two forms. \emph{Peak overshoot} is a response that rises above its direct impact; \emph{cumulative amplification} is a summed response much larger than direct impact. The data support the second form, not the first.

    \subsection{A book with depletion and refill}
Let $x_t$ be the price displacement, $d_t$ the signed depth imbalance (bid minus ask depth near the mid, relative to total) and $u_t$ signed order flow. Flow moves the price directly and depletes one side of the book, which refills at rate $r$; depth imbalance moves the price:
\begin{align}
x_t &= a\,x_{t-1} + \varphi\,d_{t-1} + \kappa\,u_t, \label{p2:eq:x}\\
d_t &= (1-r)\,d_{t-1} + \varepsilon\,u_t, \label{p2:eq:d}
\end{align}
with $\kappa>0$, $a\in[0,1)$ and $r\in(0,1]$. Write $b=1-r$. With the sign convention $d=(B-A)/(B+A)$, a buy that consumes ask depth raises $d$, so $\varepsilon>0$ means \emph{depletion} of the side that is hit, and $\varphi>0$ means that a bid-heavy book pushes the price up. The impulse response to a unit flow shock at $t=0$ is $h_0=\kappa$ and
\begin{equation}\label{p2:eq:rec}
h_k=a\,h_{k-1}+\varphi\varepsilon\,b^{k-1},\qquad k\ge1.
\end{equation}

\begin{definition}
The book \emph{amplifies} if $\max_kh_k>h_0$: the price keeps moving in the direction of the shock after the shock has passed.
\end{definition}

    \subsection{Peak overshoot}
Add momentum traders whose flow $f_t=(1-\eta)f_{t-1}+\eta q\,x_{t-1}$ enters both \eqref{p2:eq:x} and \eqref{p2:eq:d} alongside $u_t$; the two-state model is the case $q=0$. Write $C_2\equiv\varphi\varepsilon-(1-a)\kappa$ and $C_3\equiv\varphi\varepsilon+\kappa^2\eta q-(1-a)\kappa$.

\begin{theorem}[Peak overshoot: depletion versus momentum]\label{p2:thm:peak}
\begin{enumerate}[label=(\roman*)]
\item\label{p2:thm:t6} \emph{Two-state condition.} Without momentum ($q=0$) the book amplifies if and only if $C_2>0$; when $C_2\le0$ the response peaks at impact, $\max_kh_k=\kappa$.
\item\label{p2:prop:shape} \emph{Shape.} Let $q=0$, $0<a<1$ and $0<b<1$. If $a\neq b$, $h_k=Aa^k+Bb^k$ with $B=\varphi\varepsilon/(b-a)$ and $A=\kappa-B$; if $a=b$, $h_k=a^{k-1}(a\kappa+\varphi\varepsilon k)$ for $k\ge1$. When the book amplifies, $h$ viewed as a function of real $k\ge0$ has a single maximum, at
\[
k^*=\frac{\log\bigl(-A\log a/(B\log b)\bigr)}{\log(b/a)}\quad(a\neq b),\qquad k^*=\frac1{\log(1/a)}-\frac{a\kappa}{\varphi\varepsilon}\quad(a=b),
\]
the second being the limit of the first as $b\to a$; the peak lag is $\lfloor k^*\rfloor$ or $\lceil k^*\rceil$ (and is $1$ when $k^*<1$, because $h_1>h_0$). The cumulative response is $\sum_kh_k=\kappa/(1-a)+\varphi\varepsilon/[(1-a)r]$ in all cases, and faster refill lowers both the peak and the cumulative response. When $a=0$ ($h_k=\varphi\varepsilon b^{k-1}$) or $r=1$ ($h_k=a^{k-1}(a\kappa+\varphi\varepsilon)$) the response peaks at lag 1.
\item\label{p2:prop:t6m} \emph{One-period margin with momentum.} $h_1-h_0=C_3$, so $C_3>0$ implies that the book amplifies.
\item\label{p2:thm:nomom} \emph{Momentum alone cannot amplify.} If $\varphi\varepsilon=0$, $q\ge0$ and the system is stable, then $h_k=c_1\lambda_1^k+c_2\lambda_2^k$ with residues $c_1,c_2\ge0$, $c_1+c_2=\kappa$ and $|\lambda_2|\le\lambda_1=\rho<1$, so $h_k\le\kappa\rho^k<\kappa$ for every $k\ge1$; in particular $C_3<0$.
\item\label{p2:thm:noover} \emph{A sufficient condition for no overshoot.} For any signs of $\varphi$, $\varepsilon$ and $q$, if
\[
\bar C_3\equiv|\varphi\varepsilon|\Bigl(1+\frac{\kappa|q|}{r}\Bigr)+\kappa^2|q|-(1-a)\kappa\;\le\;0,
\]
then $\sup_k|h_k|=\kappa$. When $\varphi\varepsilon\ge0$ and $q\ge0$, $C_3>0\Rightarrow$ amplification $\Rightarrow\bar C_3>0$.
\end{enumerate}
\end{theorem}

\begin{proof}
(i) \emph{If.} $h_1-h_0=a\kappa+\varphi\varepsilon-\kappa=C_2>0$. \emph{Only if.} Suppose $C_2\le0$. If $\varphi\varepsilon<0$, then by \eqref{p2:eq:rec} $h_k=a^k\kappa+\varphi\varepsilon\sum_{j<k}a^{k-1-j}b^j\le a^k\kappa\le\kappa$. If $\varphi\varepsilon\ge0$, let $\Delta_k=h_k-h_{k-1}$. Then $\Delta_1=C_2\le0$ and, from \eqref{p2:eq:rec}, $\Delta_{k+1}=a\Delta_k+\varphi\varepsilon(b^{k}-b^{k-1})=a\Delta_k-\varphi\varepsilon\,r\,b^{k-1}\le a\Delta_k$, so by induction $\Delta_k\le0$ for all $k$ and $h$ peaks at $h_0=\kappa$.

(ii) Solving \eqref{p2:eq:rec} gives the two forms. For $a\ne b$, $\frac{d}{dk}h=Aa^k\log a+Bb^k\log b$ has at most one zero because $(b/a)^k$ is monotone, which gives $k^*$; for $a=b$, $\frac{d}{dk}a^{k-1}(a\kappa+\varphi\varepsilon k)=a^{k-1}[(a\kappa+\varphi\varepsilon k)\log a+\varphi\varepsilon]$ vanishes once, at the stated $k^*$, and expanding the first formula in $b-a$ gives the same limit. A unimodal function of real $k$ attains its integer maximum at $\lfloor k^*\rfloor$ or $\lceil k^*\rceil$. Summing the geometric series gives the cumulative response, which falls in $r$; the peak falls in $r$ because $h_k$ is decreasing in $r$ term by term for $\varphi\varepsilon>0$.

(iii) With a unit shock at $t=0$, $x_0=\kappa$, $d_0=\varepsilon$ and $f_1=\eta q\kappa$, so $h_1=a\kappa+\varphi\varepsilon+\kappa^2\eta q$.

(iv) Without the depth channel the state $(x_t,f_t)$ evolves by $M=\begin{psmallmatrix}a+\kappa\eta q&\kappa(1-\eta)\\ \eta q&1-\eta\end{psmallmatrix}$, which is entrywise nonnegative, and $h_k=\kappa(M^k)_{11}$. The off-diagonal product is nonnegative, so the eigenvalues are real, and by the Perron--Frobenius theorem $\lambda_2\le m_{11}\le\lambda_1=\rho$. Writing $(M^k)_{11}=c_1'\lambda_1^k+c_2'\lambda_2^k$ with $c_1'=(m_{11}-\lambda_2)/(\lambda_1-\lambda_2)\in[0,1]$ and $c_2'=1-c_1'$ gives the residues $c_j=\kappa c_j'\ge0$ and $(M^k)_{11}\le\rho^k$. Since $m_{11}\le\rho<1$, $C_3=\kappa(m_{11}-1)<0$.

(v) Suppose $|x_j|\le\kappa$ for $j<k$. Then $|f_j|\le|q|\kappa$ for $j\le k$ and $|d_{k-1}|\le|\varepsilon|\bigl(b^{k-1}+\sum_{j\ge1}b^{k-1-j}|q|\kappa\bigr)\le|\varepsilon|(1+|q|\kappa/r)$, so $|x_k|\le a\kappa+|\varphi\varepsilon|(1+\kappa|q|/r)+\kappa^2|q|\le\kappa$. Induction from $x_0=\kappa$ completes the proof.
\end{proof}

\paragraph{Interpretation.} The left side $\varphi\varepsilon$ of $C_2$ is the extra price pressure that a unit of flow creates one period later through the book it leaves behind. The right side $(1-a)\kappa$ is the part of the direct impact that reverts in one period. The book amplifies exactly when the depth channel adds more than mean reversion removes. Refill $r$ does not enter the condition; it governs how long an overshoot lasts.

    The converse of part~\ref{p2:prop:t6m} fails with two slow modes: near the stability boundary, or with slow refill and slow momentum, a book can have $C_3<0$ and still peak far above its impact, because the depth channel creates a negative residue on an intermediate mode. No condition on the one-period margin alone is exact.

    \subsection{Cumulative transient risk}
\begin{theorem}[Summed response in closed form]\label{p2:thm:sum}
Let $\tilde\kappa=\kappa+\varphi\varepsilon/r$ (long-run impact: direct plus depth) and $m=1-a$ (reversion). In the stable three-state model,
\[
S\equiv\sum_{k\ge0}h_k=\frac{\tilde\kappa}{m-q\tilde\kappa},
\]
and when $\varphi\varepsilon\ge0$ and $q\ge0$ the system is stable only if $m>q\tilde\kappa$. When all $h_k\ge0$, $S$ is the summed absolute response that enters the guard of Section~\ref{p2:sec:guard}.
\end{theorem}

\begin{proof}
$S=e_1^\top(I-M_3)^{-1}(\kappa,\varepsilon,0)^\top$ with $M_3$ the transition matrix of $(x,d,f)$. Direct computation gives $\det(I-M_3)=\eta\,[\,r(m-\kappa q)-\varphi\varepsilon q\,]$ and the stated ratio; the formula matches the simulated sums to a relative error of $7\times10^{-14}$ in 3{,}000 stable draws. If $\varphi\varepsilon\ge0$ and $q\ge0$, $M_3$ is entrywise nonnegative and stability requires $\det(I-M_3)>0$, i.e.\ $m>q\tilde\kappa$.
\end{proof}

Cumulative transient risk is long-run impact divided by reversion net of momentum. Two practical consequences follow. First, the statistic that matters for losses is the summed response, which grows without bound as reversion falls to $q\tilde\kappa$, even when no single lag overshoots (Theorem~\ref{p2:thm:peak}\ref{p2:thm:nomom}). Second, $C_3$ is best read as a one-period amplification margin: it moves with $m$, $q$ and $\tilde\kappa$ in the same direction as $S$ and is what we estimate in Sections~\ref{p2:sec:evidence} and~\ref{p2:sec:thin}. Both statistics are invariant in sign to rescaling flow.

    \section{Identification and financial safety}\label{p2:sec:infosafety}

    Chapter~\ref{ch:p1} shows that identifying the cohorts from prices is expensive. A pre-trade control does not need them; it needs a bound on the risk they create. This section asks what information certifies a safe decision, what it costs, and how much safe trading is lost without it.

    \subsection{Observation and risk contracts}\label{p2:sec:contracts}
We work in the positive cohort economy \eqref{p2:eq:J} with $q_i\ge0$ and known $a$ and $\kappa$. Its merged measure $M=\sum_iq_i\eta_i\delta_{r_i}$, $r_i=1-\eta_i$, has kernel $k_m=\int r^m\,dM$, and the price response to a unit of flow obeys $h_0=\kappa$, $h_{t+1}=ah_t+\kappa\sum_{m\le t}k_mh_{t-m}$. All $h_t\ge0$ (Chapter~\ref{ch:p1}, positive-geometry lemma), and the summed response is
\begin{equation}\label{p2:eq:SPhi}
S=\sum_th_t=\frac{\kappa}{1-a-\kappa\Phi},\qquad \Phi=\int\frac{dM(r)}{1-r}=\sum_iq_i ,
\end{equation}
with stability exactly when $a+\kappa\Phi<1$. A \emph{memory bound} $\bar r<1$ states that every memory rate satisfies $r_i\le\bar r$.

\emph{Risk contract.} A position $I$ is held fixed, each period's flow shock is bounded by $\bar u$, losses are marked to market at every date, and the horizon is unbounded. By Lemma~\ref{p2:thm:guard} the worst-case loss is $|I|\bar uS$, so the safe capacity, the largest shock bound with loss at most the budget $\mathcal B$, is $\mathcal B/(|I|S)$. Section~\ref{p2:sec:guard} adds admission, margin and fills. All statements below are deterministic worst cases under this contract; none is a probability-of-ruin statement.

\emph{Observation contracts.} We consider four: (O1) $N$ independent reset experiments, each returning $h_0,\ldots,h_{T-1}$ with independent Gaussian errors of standard deviation $\sigma$ per tap, the controlled-probe contract of Chapter~\ref{ch:p1}; (O2) an exact window $h_0,\ldots,h_n$; (O3) class-tagged flows; and (O4) audited exposure, which reveals $\Phi$. The theorems concern these contracts. Observational public records of a live book are a different contract, for which we offer evidence (Section~\ref{p2:sec:evidence}) but no finite-sample guarantee.

Given data $D$, let $\Theta(D)$ be the set of economies compatible with it. A decision is \emph{certifiably safe} if its worst-case loss is within budget for every economy in $\Theta(D)$, \emph{certifiably unsafe} if it exceeds the budget for every one, and \emph{undecidable} otherwise. This classification is a definition, not a result; it is the finance version of the data-informativity view of control without identification \citep{p2:vanwaarde2020} and of set-membership estimation \citep{p2:milanese1991}. The results below make it explicit for this economy.

    \subsection{Two markets that look the same}
\begin{example}[An exact ambiguity witness]\label{p2:ex:witness}
Let $a=0.4$ and $\kappa=1$. Market A has one cohort ($q=0.08$, $\eta=0.5$); market B has two ($q=0.02/0.95$, $\eta=0.95$; $q=0.4$, $\eta=0.05$). Their merged measures have the same mass and first moment, so in rational arithmetic their price responses coincide at lags 0, 1 and 2 ($1$, $11/25$, $267/1250$) and first differ at lag 3. Their summed responses are $1.92$ and $5.59$, so for a unit position and budget the safe shock bound is $0.52$ in A and $0.18$ in B.
\end{example}

\begin{proposition}[Unavoidable loss of safe capacity]\label{p2:prop:twoworld}
Let two economies have safe capacities $v_0<v_1$ and data laws $P_0,P_1$. If a data-driven rule $A(D)\in[0,v_1]$ admits an unsafe level in economy 0 with probability at most $\beta$, $P_0(A>v_0)\le\beta$, then
\[
\E_1[v_1-A(D)]\;\ge\;(v_1-v_0)\,\bigl[1-\beta-{\rm TV}(P_0,P_1)\bigr]_+ .
\]
The bound is attained when $P_0=P_1$ by the rule that admits $v_1$ with probability $\beta$ and $v_0$ otherwise.
\end{proposition}

\begin{proof}
Let $E=\{A\le v_0\}$. Then $P_0(E)\ge1-\beta$ and $P_1(E)\ge P_0(E)-{\rm TV}$. On $E$ the capacity given up in economy 1 is at least $v_1-v_0$, and it is nonnegative elsewhere.
\end{proof}

The argument is the two-point method of statistical decision theory \citep{p2:lecam1986}; its use here is to price missing information in units of safe trading. In Example~\ref{p2:ex:witness}, a supervisor who sees the first three coefficients and must admit an unsafe level in B with probability at most 5\% gives up at least $0.95\times0.34=0.32$ units of safe shock capacity in A.

    \subsection{What a finite window certifies}
Lags $1$ to $n$ of the response determine the first $n$ moments of $M$ exactly, through $k_t=(h_{t+1}-ah_t-\kappa\sum_{m<t}k_mh_{t-m})/\kappa^2$. The summed response depends on $M$ only through $\Phi=\int f\,dM$ with $f(r)=(1-r)^{-1}$. Certifying it from a window is therefore an extremal moment problem.

\begin{theorem}[Finite-window certification]\label{p2:thm:window}
Observe $h_0,\ldots,h_n$ exactly under a memory bound $\bar r<1$, and suppose at least one compatible economy is stable. Suppose the implied moment vector $(k_0,\ldots,k_{n-1})$ lies in the interior of the moment space $\mathcal M_n([0,\bar r])$ (equivalently, the Hankel matrices of the moment sequence and of its localization by $r(\bar r-r)$ are positive definite). Let $\mu_-$ and $\mu_+$ be its lower and upper principal representations on $[0,\bar r]$ (the extremal measures of the Markov--Krein theorem, unique in the interior; $\mu_+$ has an atom at $\bar r$), and $\Phi_\pm=\int f\,d\mu_\pm$. On the boundary of the moment space the compatible measure is unique and $\Phi_-=\Phi_+$.
\begin{enumerate}[label=(\roman*)]
\item \emph{Sharp interval.} The summed responses of stable economies compatible with the data form the interval $[S(\Phi_-),S(\Phi_+)]$ if $a+\kappa\Phi_+<1$, and $[S(\Phi_-),\infty)$ otherwise. Every finite endpoint in the stable class is attained by an economy with at most $\lfloor n/2\rfloor+1$ cohorts. When the upper endpoint is infinite it is a supremum approached by stable economies, not an attained response.
\item \emph{Geometric convergence.} $\Phi_+-\Phi_-\le2k_0E_{n-1}$, where $E_{n-1}=(2/\bar r)\,\varrho^{\,n-1}/(x_0^2-1)$ is the error of the best uniform approximation of $f$ on $[0,\bar r]$ by polynomials of degree $n-1$, $x_0=2/\bar r-1$ and $\varrho=x_0-\sqrt{x_0^2-1}<1$. A window of length $n=O\bigl(\log(1/\epsilon)/\log(1/\varrho)\bigr)$ certifies $\Phi$ to precision $\epsilon$, and $\log(1/\varrho)\sim2\sqrt{(1-\bar r)/\bar r}$ as $\bar r\to1$.
\item \emph{No certificate without a memory bound.} If memory rates may lie anywhere in $[0,1)$, then for every $n$ and every window whose moment vector is interior to the moment space, compatible stable economies have arbitrarily large summed responses: no finite window certifies any finite risk bound.
\end{enumerate}
\end{theorem}


    Part (i) applies the Markov--Krein theorem \citep{p2:krein1977,p2:karlin1966}, part (ii) Chebyshev's best approximation of $1/(x_0-x)$ \citep{p2:achieser1956}; the proofs are in Appendix~\ref{p2:app:proofs}. In an eight-cohort economy with memory bound $0.95$ and true summed response $5.56$, lags $0$--$3$ certify nothing, lags $0$--$6$ give $[5.01,7.20]$ and lags $0$--$8$ give $[5.30,6.00]$; for fixed data the upper end grows without limit as the memory bound approaches one, as part (iii) says. The endpoints are proved, not only computed: polynomial bounds whose sign conditions are verified in interval arithmetic with outward rounding turn them into certified intervals \citep{p2:rump2010}.

    \subsection{What certification costs, and what identification costs}
\begin{theorem}[Certification cost of the summed response]\label{p2:thm:certcost}
Consider the class $\mathcal K(\bar r,\bar s)$ of positive economies with memory bound $\bar r$ and $a+\kappa\Phi\le\bar s<1$, observed under contract (O1). Let $\bar\lambda=1-(1-\bar r)(1-\bar s)$ and $\bar h_t$ the average of the $N$ responses.
\begin{enumerate}[label=(\roman*)]
\item \emph{A valid certificate.} For every economy in the class, $h_t\le\kappa\bar\lambda^t$, and
\[
\hat S^+_T=\sum_{t<T}\bar h_t+z_{1-\gamma}\,\sigma\sqrt{T/N}+\frac{\kappa\bar\lambda^T}{1-\bar\lambda}
\]
satisfies $P(S\le\hat S^+_T)\ge1-\gamma$. With $T_\epsilon=\lceil\log(2\kappa/((1-\bar\lambda)\epsilon))/\log(1/\bar\lambda)\rceil$ and $N\ge16z_{1-\gamma}^2\sigma^2T_\epsilon/\epsilon^2$, the excess $\hat S^+-S$ is at most $\epsilon$ with probability at least $1-2\gamma$. Certification to precision $\epsilon$ therefore costs $N=O(\sigma^2\epsilon^{-2}\log(1/\epsilon))$.
\item \emph{A matching lower bound.} Fix an interior single-cohort economy $\theta_0$ and let $C_0=\|\partial h/\partial q\|_2/(\partial S/\partial q)$ at $\theta_0$. Any procedure that is valid at level $1-\gamma$ on the class and overshoots $S$ by less than $2\epsilon$ with probability at least $1/2$ at $\theta_0$ needs $N\ge\bigl((1-2\gamma)\sigma/(2C_0\epsilon)\bigr)^2(1+o(1))$ as $\epsilon\to0$.
\item \emph{The logarithm is necessary.} There is $c>0$ depending only on $(a,\kappa,\bar r,\bar s)$ such that any procedure that is valid at level $1-\gamma$, $\gamma<1/4$, on the class and overshoots $S$ by at most $\epsilon$ with probability at least $1-\gamma$ at every economy of the class needs $N\ge c\,\sigma^2\epsilon^{-2}\log(1/\epsilon)$ for all small $\epsilon$, however many responses each experiment records. With part (i), the cost of certifying $S$ is $\Theta(\sigma^2\epsilon^{-2}\log(1/\epsilon))$.
\item \emph{Independence of the cohort structure.} None of these bounds depends on the number of cohorts or on the spacing of their memories, whereas identifying them costs $N\asymp\sigma^2(\Delta\eta)^{-(4K-2)}$ with exponential memory and $\exp(\Theta(1/w))$ experiments for accuracy $w$ with general memory (Chapter~\ref{ch:p1}).
\end{enumerate}
\end{theorem}

The certificate in part (i) is elementary: it is a sum of averaged taps, a Gaussian margin and a tail that the memory bound controls. Part (iii) shows that its logarithm is not an artefact of truncation: no procedure, however it weights the observed responses, certifies the summed response at the parametric rate, because $1/(1-r)$, the kernel of the total feedback, has a pole at infinite memory and any polynomial that approximates it on the memory range must carry large coefficients. Part (ii) is the parametric bound at a fixed economy; part (iii) is the worst case over the class. The content is in the comparison. The compared inverse problems differ: \citet{p2:tu2018} study $\mathcal H_\infty$-norm estimation relative to FIR input-output model identification, whereas the comparison here concerns a positive response functional relative to latent cohort recovery. These are distinct tasks; no general claim that norm certification is easier than identifying the observable response follows. Table~\ref{p2:tab:certvsid} makes this concrete on Chapter~\ref{ch:p1}'s hardest two-cohort pairs. Those pairs have \emph{identical} summed responses, because they share the total weight $\sum_iq_i$: the directions in which prices are least informative are directions in which the risk a guard must bound does not move. Identifying their weights to $\pm0.01$ needs from 7 to 1.5 million reset experiments as the spacing falls from $0.30$ to $0.04$ (log--log slope $-6.06$), while certifying the summed response to 1\% needs 11 to 18 at every spacing (slope $+0.23$).

\begin{table}[t]\centering\small
\begin{threeparttable}
\caption{Identifying the cohorts versus certifying the risk they create}\label{p2:tab:certvsid}
\begin{tabular}{lcccccccc}
\toprule
Rate spacing $\Delta\eta$ & 0.30 & 0.20 & 0.14 & 0.10 & 0.08 & 0.06 & 0.05 & 0.04\\
\midrule
Identify weights ($\pm0.01$), $N$ & 7 & 88 & 775 & 5{,}928 & 22{,}741 & 128{,}333 & 383{,}859 & 1{,}466{,}362\\
Certify $S$ (1\%), $N$ & 18 & 14 & 13 & 12 & 11 & 11 & 11 & 11\\
\bottomrule
\end{tabular}
\begin{tablenotes}\footnotesize\item The two-cohort moment-matched pairs of Chapter~\ref{ch:p1} ($q=0.2$, mean rate $0.4$, $\tau=0.01$, $a=0.5$, $\kappa=0.8$, per-tap noise $1.6\times10^{-3}$). Identification: Le Cam requirement at 95\% (Chapter~\ref{ch:p1}, Table on two-point sample complexity). Certification: sufficient $N$ from Theorem~\ref{p2:thm:certcost}(i) with memory bound equal to the largest rate in the pair and $\bar s=a+\kappa q+0.02$. The two economies of each pair have the same summed response, $2.353$. Kind: recomputed (\texttt{research\_extension/information\_safety}, test T03).\end{tablenotes}
\end{threeparttable}
\end{table}

    Two refinements are in the full version. If the noise in the taps of one experiment is correlated, the certificate of part (i) can under-cover (75\% at nominal 95\% in simulation), while a Student-$t$ certificate on the summed taps and a Hotelling region over the window keep their level \citep{p2:hotelling1931}; and experiments can be run one at a time with an anytime-valid version of the certificate that stops once the data decide.

    \section{An identification-free pre-trade guard}\label{p2:sec:guard}

    \subsection{Maximality}
A dealer with inventory $I$ chooses an admission level $\alpha\in[0,1]$ (the fraction of incoming algorithmic flow it accepts, or the depth it quotes). The closed-loop operator is $J(\alpha)=J_0+\alpha B_F$ with price impulse response $g_k(\alpha)$. Shocks satisfy $|u_t|\le\bar u$. Over horizon $H$ under a reset initial state $z_0=0$ (or after subtracting a known or directly observed initial drift $b_t=CJ(\alpha)^t z_0$ from the remaining budget), the shock-driven price change is $\sum_{k<H}g_k(\alpha)u_{t-k}$, and the dealer has a loss budget $\mathcal B$. The contract is: inventory $I(\alpha)$ is held fixed over the horizon, the initial state is reset ($z_0=0$) or its drift is observed, each period's shock is bounded by $\bar u$, and loss is the mark-to-market loss at any date in the horizon; because the partial sums $\sum_{k<h}|g_k|$ increase in $h$, the worst date is the last. Margin, hedging, fills and non-zero initial states are treated in the full version.

\begin{lemma}[Summed-gain guard is maximal]\label{p2:thm:guard}
The worst-case loss over horizon $H$ is
\[
L_{\max}(\alpha)=|I(\alpha)|\,\bar u\sum_{k<H}|g_k(\alpha)|.
\]
Hence the guard ``admit $\alpha$ iff $L_{\max}(\alpha)\le\mathcal B$'' is necessary and sufficient for bounded-shock solvency: it admits every safe $\alpha$ and no unsafe one. Replacing $\sum_k|g_k|$ by $\max_k|g_k|$ can admit unsafe levels; replacing it by $H\max_k|g_k|$ is safe but conservative.
\end{lemma}

\begin{proof}
By H\"older's inequality $|\sum_kg_ku_{t-k}|\le\bar u\sum_k|g_k|$, with equality for $u_{t-k}=\bar u\,{\rm sign}(g_k)$, which is admissible. Multiplying by $|I(\alpha)|$ gives $L_{\max}$, so the admission set equals the safe set. Since $\max_k|g_k|\le\sum_k|g_k|\le H\max_k|g_k|$, the peak guard under-states loss and the $H$-scaled guard over-states it.
\end{proof}

    The identity is the classical $\ell_1$ gain of a convolution operator \citep{p2:vidyasagar1986,p2:dahleh1987,p2:dahleh1995}; we use it rather than claim it. The finance content is that the dealer's admission level changes the operator, so the guard is a rule over an action set; that its single input is invariant to how the traders split into cohorts and is identified from prices; and that the same statistic prices the corrective rebate in equilibrium (Proposition~\ref{p2:prop:equil}). The guard coincides with a brute-force adversary's maximal safe level in 200 of 200 random states, and the peak-gain guard admitted an unsafe level in 2{,}000 of 2{,}000. An infinite-horizon gain is computed on a finite horizon with a contraction-block tail bound.

    \subsection{Robustness to estimation error}
\begin{theorem}[Robust gain]\label{p2:thm:robust}
Let $M=\sup_k\|J^k\|$ and $S=\sum_k\|J^k\|$. If $\|E\|\le\epsilon$ and $\epsilon S<1$, then $\sup_k\|(J+E)^k\|\le M/(1-\epsilon S)$.
\end{theorem}

\begin{proof}
Let $P_k=(J+E)^k$. Expanding, $P_k=J^k+\sum_{j<k}J^{k-1-j}EP_j$. With $m_k=\|P_k\|$, $m_k\le M+\epsilon\sum_{j<k}\|J^{k-1-j}\|m_j\le M+\epsilon S\max_{j<k}m_j$. Since $m_0=1\le M$, induction gives $m_k\le M/(1-\epsilon S)$ for all $k$.
\end{proof}

Theorem~\ref{p2:thm:robust} is deterministic: it holds whenever $\|E\|\le\epsilon$. A statistical guarantee follows only through a confidence set for the operator, and then holds with that set's coverage probability, not with certainty. The bound held in all 30,000 random perturbations of 100 non-normal operators, with median tightness $0.50$ (realized gain over bound; maximum $0.57$). In a simulated window estimator of the kind used in Section~\ref{p2:sec:evidence} (600 transitions), a residual-bootstrap 95\% ball for the operator covered the true $J$ in $94.7\%$ of 300 replications, with median radius $0.086$; but $\epsilon S<1$ held in only $8\%$ of windows. The certificate is valid but often unavailable at realistic sample sizes; longer windows or pooled estimates are needed to use it in practice.

The guard needs the summed gain, which is invariant to cohort splitting (Proposition~\ref{p2:prop:invariance}) and certified with $\Theta(\sigma^2\epsilon^{-2}\log(1/\epsilon))$ experiments whatever the spacing of the memories (Theorem~\ref{p2:thm:certcost}); it does not need the algorithms' identities, which are expensive (Chapter~\ref{ch:p1}). In a two-cohort economy with rates $0.7\pm\Delta$, the relative error of the estimated summed gain is flat in $\Delta$ (log--log slope $-0.15$), while the error of the cohort weights grows as $\Delta^{-3.07}$.

    \subsection{An uncertainty-aware guard}\label{p2:sec:robustguard}
The guard of Lemma~\ref{p2:thm:guard} assumes the operator is known. With data, the relevant question is which admission levels are safe for every economy the data allow. We answer it for a venue gate (a venue that admits a fraction of incoming algorithmic flow) in the positive cohort economy, where admitting a fraction $\alpha$ of the algorithmic flow scales every cohort's weight, $M\mapsto\alpha M$, so that $S(\alpha)=\kappa/(1-a-\kappa\alpha\Phi)$.

\begin{lemma}[Soundness of compatible-set guards]\label{p2:lem:sound}
If a data-dependent set $\mathcal C(D)$ contains the true economy with probability at least $1-\gamma$, then every action whose worst-case loss is within budget for all economies in $\mathcal C(D)$ is within budget for the true economy with probability at least $1-\gamma$, including an action selected from the data. If no action qualifies, the guard certifies nothing: it halts new admission and flags the state, because existing inventory can be unsafe even with no new trading.
\end{lemma}

\begin{theorem}[Optimal uncertainty-aware admission]\label{p2:thm:robustguard}
Under the reset or observed-drift contract ($z_0=0$, or known drift subtracted from $\mathcal B$), let position $I$, shock bound $\bar u$ and budget $\mathcal B$ be given, and let $\Phi^+$ be an upper bound on $\Phi$ valid with probability at least $1-\gamma$ (for example $\Phi^+=(1-a-\kappa/\hat S^+)/\kappa$ from Theorem~\ref{p2:thm:certcost}, or $\Phi_+$ from Theorem~\ref{p2:thm:window}). Put $c_I=1-a-\kappa|I|\bar u/\mathcal B$.
\begin{enumerate}[label=(\roman*)]
\item \emph{Closed form and interval.} If $c_I\ge0$, the levels that are safe for every $\Phi\le\Phi^+$ form the interval $[0,\alpha^*]$ with $\alpha^*=\min\{1,\ c_I/(\kappa\Phi^+)\}$.
\item \emph{Safety.} With probability at least $1-\gamma$, every admitted level is safe for the true economy.
\item \emph{Maximality.} Every level above $\alpha^*$ is unsafe for an economy with $\Phi=\Phi^+$; when $\Phi^+$ is attained by a compatible economy (as for $\Phi_+$ in Theorem~\ref{p2:thm:window}), no rule using the same information can admit more without admitting an unsafe level.
\item \emph{Infeasibility.} If $c_I<0$, no level is safe, including $\alpha=0$: the guard halts admission and the position must fall to $|I|\le\mathcal B(1-a)/(\kappa\bar u)$.
\item \emph{Minimax capacity loss.} On the class of Theorem~\ref{p2:thm:certcost} with $\Phi$ bounded away from zero, the certificate-based guard gives up $\alpha^\circ-\alpha^*=O\bigl(\sigma\sqrt{\log N/N}\bigr)$ of the oracle capacity $\alpha^\circ=\min\{1,c_I/(\kappa\Phi)\}$ with probability at least $1-2\gamma$, and every rule that admits an unsafe level with probability at most $\beta<1$ gives up at least $c\,\sigma/\sqrt N$ in expectation at some economy of the class. The guard is minimax-optimal in lost capacity up to a logarithmic factor.
\end{enumerate}
\end{theorem}

Monotonicity of $S(\alpha)$ in $\alpha$ and $\Phi$ gives (i)--(iv); (v) combines the certificate's excess with Proposition~\ref{p2:prop:twoworld} applied to two economies whose capacities differ by order $\sigma/\sqrt N$ (Appendix~\ref{p2:app:proofs}). In 500 random economies the closed form matched a brute-force search over 1{,}001 admission levels, computed from the state matrices rather than from \eqref{p2:eq:SPhi}, in every case; every safe set was an interval; and the 28 infeasible states were flagged (test T06). Across 18 designs the certificate covered the true summed response at least at its nominal level (minimum 93.6\% at nominal 90\%), with a median excess of 3--15\% when the memory bound is $0.6$ or $0.9$ and about 230\% when it is $0.97$, where the tail term dominates (test T04).

    When admission also changes the book, deepening it and speeding reversion as well as adding feedback, risk can first rise and then fall with the admission level, and the safe levels need not form an interval: certifying the largest safe level then no longer certifies those below it. The full version certifies the whole safe set by interval branch and bound \citep{p2:hansen2004}, with soundness and completeness; in a book where admission raises feedback linearly and liquidity quadratically the certified safe set is $[0,0.053]\cup[0.178,1]$.

    \paragraph{Why the peak guard is unsafe in a real book.} At the median BTCUSDT estimates of Section~\ref{p2:sec:evidence}, price impact is persistent ($\hat a\approx0.97$ per minute), so the summed gain over 15 minutes is about 13 times the peak. A guard calibrated to the peak would admit inventory about 13 times larger than is safe against a sustained adversarial flow of the declared size, although that book does not amplify. Persistence alone makes peak-based limits unsafe.

    \paragraph{The guard on a dealer's balance sheet.} The full version carries the guard through maintenance margin, hedging and worst-case fills, and proves forward solvency at every step at which an admissible level exists; when none exists the guard halts new fills and flags the state, because existing inventory can still breach margin. In a matched-shock Monte Carlo inside the model, with every rule calibrated to the same trading capacity, the guard has no forced liquidations against 31.6\% with no control, 31.4\% under a trailing volatility limit and 20.3\% under an inventory cap. With a hidden cohort state the drift of future prices is not identification-invariant under a Euclidean bound, whereas the natural interval bound supports the following sufficient invariant box. Cumulative unforced drift still depends on the updating rates:
    
\begin{proposition}[Hidden-state drift and an invariant-box guard]\label{p2:prop:driftguard}
Let $0\le a<1$, $\kappa>0$, $q_i\ge0$, $0<\eta_i\le1$ and $\Phi=\sum_iq_i\le\Phi^+$. Admission $\alpha\in[0,1]$ scales cohort demand, with $a+\alpha\kappa\Phi^+<1$. Use the original price-first coordinates $x_{t+1}=ax_t+\alpha\kappa\sum_iq_iy_{i,t}+\kappa u_{t+1}$ and $y_{i,t+1}=(1-\eta_i)y_{i,t}+\eta_ix_{t+1}$.
\begin{enumerate}[label=(\roman*)]
\item \emph{Cumulative unforced drift.} For fixed weights and rates and $u=0$,
\begin{equation}\label{p2:eq:driftenv}
\sup_{\|y_0\|_\infty\le\bar y_0}\sum_{t=1}^\infty|x_t|
=\frac{a|x_0|+\alpha\kappa\bar y_0\sum_iq_i/\eta_i}{1-a-\alpha\kappa\Phi}.
\end{equation}
Aligned initial memories attain the bound. Thus the cumulative quantity is not determined by $\Phi$ and $\bar y_0$ alone when rates can approach zero.
\item \emph{Exact finite-horizon pointwise envelope.} Write $c^\top=e_x^\top$, $b=(\kappa,\kappa\eta_1,\ldots,\kappa\eta_K)^\top$ and $\bar z_0=(|x_0|,\bar y_0\mathbf1^\top)^\top$. For $|u_t|\le\bar u$, the worst-case peak over $0\le t\le H$ is
\begin{equation}\label{p2:eq:driftpeak}
\max_{0\le t\le H}\left[c^\top J(\alpha)^t\bar z_0+\bar u\sum_{j=0}^{t-1}c^\top J(\alpha)^jb\right].
\end{equation}
The empty sum at $t=0$ is zero. Aligned memories and constant shocks attain every term simultaneously. This formula retains the time index shared by initial drift and forced response.
\item \emph{Identification-free sufficient box.} For fixed inventory $I\ne0$ and budget $\mathcal B>0$, let $L=\mathcal B/|I|$. If $|x_0|\le L$, $\bar y_0\le L$ and $aL+\alpha\kappa\Phi^+L+\kappa\bar u\le L$, then $|I x_t|\le\mathcal B$ for every date and every compatible positive cohort economy. For $\Phi^+>0$, $\bar u>0$ and $(1-a)L\ge\kappa\bar u$, a sufficient admission limit is
\begin{equation}\label{p2:eq:alphabox}
\alpha_{\rm box}=\min\left\{1,\frac{(1-a)L-\kappa\bar u}{\kappa\Phi^+L}\right\}.
\end{equation}
If these initial-box or feasibility conditions fail, this argument is inconclusive. If $I=0$, the fixed-inventory mark-to-market loss is zero; if $\Phi^+=0$, admission does not affect this plant. This is not a theorem about fills, changing inventory, fees or maintenance margin.
\end{enumerate}
\end{proposition}

At the Table~\ref{p2:tab:certvsid} calibration ($a=0.4$, $\kappa=1$, $\Phi^+=0.4$, $|I|=1$, $\bar u=0.48$, $\mathcal B=1$, $x_0=0$, $\bar y_0=0.25$), the invariant box admits $\alpha_{\rm box}=0.3000$, the same as the reset guard. Its capacity reduction in this calibration is $0.0000$. Indeed $0.4+0.30\times0.4+0.48=1$; every memory update is a convex combination of values in $[-1,1]$. For any larger stable admission, an economy with $\Phi=0.4$ under constant positive shocks has limit $0.48/(0.6-0.4\alpha)>1$, proving maximality in this calibration. The earlier $0.2400$ rule was conservative here; its claimed exact $0.0600$ capacity cost is withdrawn. Original result files are retained as historical records. The corrected checks are in \texttt{recomputed/R\_P2\_step2\_theory\_repaired.json} and \texttt{studies/economic\_review\_20261029/}.

    Re-optimization by the traders does not undo the guard: when a venue throttles admission, filtering cohorts lower their adaptation rates, and in the positive cohort economy every finite-horizon summed response falls with them.

    \section{The value of disclosure for safe trading}\label{p2:sec:disclosure}

    Chapter~\ref{ch:p1} measures disclosure in samples needed to identify the cohorts. For a pre-trade control the relevant value is the safe trading capacity it unlocks.
    
\begin{theorem}[Decision value of disclosure]\label{p2:thm:disclosure}
Under the risk contract of Section~\ref{p2:sec:contracts}, compare an observation contract $O$ with a richer contract $O'$ that reveals everything $O$ reveals and more, both with exact data.
\begin{enumerate}[label=(\roman*)]
\item \emph{Monotonicity.} The compatible sets are nested, $\Theta_{O'}\subseteq\Theta_O$, so $\Phi^+_{O'}\le\Phi^+_O$ and the admitted level of Theorem~\ref{p2:thm:robustguard} satisfies $\alpha^*_{O'}\ge\alpha^*_O$.
\item \emph{Strictness.} $\alpha^*_{O'}>\alpha^*_O$ if and only if $\alpha^*_O<1$, $c_I>0$, and $O'$ excludes every compatible economy with $\Phi=\Phi^+_O$. Under the window contract that economy is the upper principal representation of Theorem~\ref{p2:thm:window}: disclosure has decision value exactly when it rules out a cohort at the memory bound with the extremal weight.
\item \emph{Long memory.} Under the window contract with memory bound $\bar r$, $\alpha^*_O\to0$ as $\bar r\to1$ for any fixed window (Theorem~\ref{p2:thm:window}(iii)), whereas audited exposure, which reveals $\Phi$, admits the oracle capacity $\alpha^\circ$. The decision value of audited disclosure therefore tends to the whole oracle capacity as memory lengthens.
\item \emph{Net value.} With value $w$ per unit of admitted capacity and disclosure cost $c(O')$, the net value $w(\alpha^*_{O'}-\alpha^*_O)-c(O')$ is zero or negative whenever $O$ already identifies $\Phi$ (a window whose moment vector lies on the boundary of the moment space) or the budget does not bind.
\item \emph{Noisy proxies.} With noisy data the certificates of two contracts are confidence bounds, not nested sets, and a richer contract can admit less. Under class-tagged flows measured with noise over $T$ periods and estimated as in Chapter~\ref{ch:p1}, class $i$'s exposure $\hat\Phi_i=\hat w_i/(1-\hat r_i)$ has delta-method standard error
\[
{\rm se}(\hat\Phi_i)=\frac{1}{(1-r_i)\sqrt T}\Bigl[V_{ww}+2\Phi_iV_{wr}+\Phi_i^2V_{rr}\Bigr]^{1/2},
\]
with $V$ the asymptotic covariance of $(\hat w_i,\hat r_i)$, so the tagged bound $\sum_i\hat\Phi_i+z_{1-\gamma}{\rm se}$ admits less than the public certificate whenever its margin exceeds $\Phi^+_O-\Phi$. Intersecting the two certificates, each at level $1-\gamma/2$, gives a bound valid at level $1-\gamma$ that admits at least as much as either certificate at level $1-\gamma/2$, which restores monotonicity.
\end{enumerate}
\end{theorem}

The theorem is set inclusion plus the closed form of Theorem~\ref{p2:thm:robustguard}; its content is the location of the value. Disclosure is worth most where public certificates are weakest, in markets whose participants have long memories. In simulation (test T07; Table~\ref{p2:tab:discsim}), with two classes and one of them slow, audited exposure adds 0.027 of admitted capacity to a public certificate when the slow class's memory rate is 0.6 and 0.122 when it is 0.995, out of an oracle capacity of 0.667. Class-tagged flows measured with 10\% noise add 0.02--0.03 at memory rates up to 0.97 but \emph{lower} capacity at 0.995, because a class's exposure $q=w/(1-r)$ is badly conditioned when $r$ is near one: disclosure of a noisy proxy can be worth less than nothing. Part (v) quantifies this. In a second simulation of the same design (test T15), the standard error of the slow class's exposure rises from $0.0004$ to $0.035$ as $1-r$ falls from $0.4$ to $0.005$, roughly in proportion to $(1-r)^{-1}$ over the last decade, while the public certificate's margin grows only with the square root of its horizon. At $r=0.995$ the tagged bound admits $0.506$ against $0.549$ for the public one; the intersection admits $0.537$, close to the better of the two, and covers the true exposure in 99\% of replications (Table~\ref{p2:tab:discsim}). Chapter~\ref{ch:p1}'s evidence that Binance's and the CFTC's class positions pin down memories that prices leave unrestricted concerns class-level positions, which are proxies, not records of individual algorithms.

\begin{table}[t]\centering\small
\begin{threeparttable}
\caption{Admitted capacity under three observation contracts (simulation)}\label{p2:tab:discsim}
\begin{tabular}{lcccc}
\toprule
Memory rate of the slow class & 0.60 & 0.90 & 0.97 & 0.995\\
\midrule
Public responses (certificate, Theorem~\ref{p2:thm:certcost}) & 0.640 & 0.623 & 0.601 & 0.544\\
Class-tagged flows, 10\% noise & 0.662 & 0.655 & 0.631 & 0.507\\
Audited exposure (oracle) & 0.667 & 0.667 & 0.667 & 0.667\\
Public $\cap$ tagged, each at 97.5\% (test T15) & 0.661 & 0.653 & 0.626 & 0.537\\
\midrule
Value of audited exposure over public & 0.027 & 0.044 & 0.066 & 0.122\\
\bottomrule
\end{tabular}
\begin{tablenotes}\footnotesize\item Two classes ($q=0.06$ with the stated memory rate, $q=0.12$ with rate $0.5$), $a=0.4$, $\kappa=1$, position 1, shock bound $0.48$, budget 1; public: 500 reset experiments, per-tap noise $0.05$, certificate at 95\% with the true memory bound; tagged: 2{,}000 periods of each class's flow with 10\% measurement noise, two-stage least squares as in Chapter~\ref{ch:p1}, 95\% delta-method bound; 300 replications. Entries are mean admitted fractions $\alpha^*$. The intersection row comes from an independent run of the same design (seed 15), whose public and tagged rows are $0.639/0.662$, $0.625/0.655$, $0.600/0.631$ and $0.549/0.506$. Kind: recomputed (tests T07 and T15).\end{tablenotes}
\end{threeparttable}
\end{table}

\paragraph{Pricing versus guarding.} A conservative bound is the right input for a guard and the wrong input for a price. The corrective rebate of Proposition~\ref{p2:prop:equil} depends on $S^2$; computed from an upper bound that overstates $S$ by 10\% it over-subsidizes reversion and loses 2.4--2.8\% of first-best welfare, and by 50\%, 70--80\% (test T12). Guards should use certified upper bounds; rebates should use point estimates.

    \section{Who provides reversion? An equilibrium}\label{p2:sec:equil}

    Theorem~\ref{p2:thm:sum} says that transient risk is governed by $S=\tilde\kappa/(m-q\tilde\kappa)$, and its ingredients are chosen by participants: dealers decide how hard to lean against displacements, which sets reversion $m$, and momentum traders how strongly to chase them, which sets $q$.
    
\subsection{Setup}
There are $N$ dealers. Dealer $i$ supplies reversion intensity $s_i\ge0$ by absorbing order flow against the displacement, so $m=m_0+\sum_is_i$, where $m_0$ is reversion supplied by other means. Leaning earns a spread $v$ per unit of intensity at quadratic inventory cost, so the private value is $vs_i-\tfrac c2s_i^2$. By Lemma~\ref{p2:thm:guard} the worst-case transient loss of a position of size $Q$ facing flow shocks bounded by $U$ is $QU\,S$, so we take aggregate expected tail loss to be $\Lambda S(m,q)$ with $\Lambda>0$. Dealers bear a share $\theta<N$ of it, so each bears $\theta/N$; leveraged users and their counterparties bear the rest. Momentum traders choose aggregate intensity $q$ with private value $bq-\tfrac e2q^2$ and bear none of the tail loss. Long-run impact $\tilde\kappa$ is taken as given.

\subsection{Results}

\begin{proposition}[Equilibrium transient risk]\label{p2:prop:equil}
Let $g(m)\equiv-\partial S/\partial m=\tilde\kappa/(m-q\tilde\kappa)^2$.
\begin{enumerate}[label=(\roman*)]
\item \emph{Under-provision of reversion.} The symmetric Nash equilibrium solves $m-m_0=N v/c+\theta\Lambda g(m)/c$, the planner's optimum solves $m-m_0=N v/c+N\Lambda g(m)/c$, and $m^{\rm eq}<m^{\rm SP}$, so $S^{\rm eq}>S^{\rm SP}$.
\item \emph{Fewer dealers, more risk.} $\partial m^{\rm eq}/\partial N=(v/c)/(1-\theta\Lambda g'(m)/c)>0$, so equilibrium reversion rises and $S^{\rm eq}$ falls with the number of dealers.
\item \emph{Over-provision of momentum.} Momentum traders choose $q=b/e$; the planner chooses $q^{\rm SP}<b/e$, solving $b-eq=\Lambda\tilde\kappa^2/(m-q\tilde\kappa)^2$.
\item \emph{The corrective price is identified.} The first best for reversion is implemented by a rebate per unit of reversion
\[
\tau^*=\Bigl(1-\frac\theta N\Bigr)\Lambda\,g(m)=\Bigl(1-\frac\theta N\Bigr)\Lambda\,\frac{S^2}{\tilde\kappa},
\]
which depends on the market only through $S$ and $\tilde\kappa$.
\item \emph{Endogenous depth.} Suppose in addition that each dealer posts depth $D_i\ge0$ at private value $v_DD_i-\tfrac{c_D}2D_i^2$, that long-run impact is $\tilde\kappa=1/D$ with $D=\sum_iD_i$, and that momentum traders size their flow to the book, $q=\psi D$ with $\psi$ fixed. Then $S=\Gamma/D$ with $\Gamma\equiv S/\tilde\kappa=1/(m-\psi)$, which does not depend on depth. Symmetric equilibrium depth solves $D=Nv_D/c_D+\theta\Lambda\Gamma/(c_DD^2)$ and the planner's $D=Nv_D/c_D+N\Lambda\Gamma/(c_DD^2)$, so $D^{\rm eq}<D^{\rm SP}$; $D^{\rm eq}$ increases with $N$, so $S^{\rm eq}=\Gamma/D^{\rm eq}$ falls with the number of dealers while $S/\tilde\kappa$ does not move. The rebate per unit of depth $\tau_D^*=(1-\theta/N)\Lambda S\tilde\kappa$ implements the first-best depth.
\end{enumerate}
\end{proposition}

\begin{proof}
(i) Dealer $i$'s first-order condition is $v-cs_i+(\theta/N)\Lambda g(m)=0$; the planner's is $v-cs+\Lambda g(m)=0$. Summing over dealers gives the two fixed-point equations. The right-hand sides are decreasing in $m$ because $g$ is decreasing on $m>q\tilde\kappa$, and the planner's lies above the dealers' for every $m$ when $\theta<N$, so its crossing with the line $m-m_0$ lies to the right. (ii) follows by implicit differentiation, since $g'<0$. (iii) $\partial S/\partial q=\tilde\kappa^2/(m-q\tilde\kappa)^2>0$ is the marginal external cost of momentum that traders ignore. (iv) The rebate closes the wedge between the two first-order conditions, and $g=S^2/\tilde\kappa$ by Theorem~\ref{p2:thm:sum}. (v) With $q\tilde\kappa=\psi$, Theorem~\ref{p2:thm:sum} gives $S=\tilde\kappa/(m-\psi)=\Gamma/D$ and $\partial S/\partial D_i=-\Gamma/D^2$. Dealer $i$'s first-order condition for depth is $v_D-c_DD_i+(\theta/N)\Lambda\Gamma/D^2=0$ and the planner's $v_D-c_DD_i+\Lambda\Gamma/D^2=0$; summing gives the two fixed points. In each, the left side $D$ increases and the right side decreases in $D$, so each has a unique positive solution, the planner's to the right of the dealers' because $\theta<N$; and $D-\theta\Lambda\Gamma/(c_DD^2)=Nv_D/c_D$ has a left side increasing in $D$, so $D^{\rm eq}$ increases with $N$. The rebate closes the wedge $(1-\theta/N)\Lambda\Gamma/D^2=(1-\theta/N)\Lambda S\tilde\kappa$.
\end{proof}

At $v=0.002$, $c=0.5$, $\theta=0.3$, $\Lambda=10^{-4}$, $m_0=0.005$, $q=0.02$ and $\tilde\kappa=1$, equilibrium $S$ falls from 27.9 with one dealer to 6.8 with forty, and the planner's $S$ is 1.5--2.7 times lower throughout; the rebate computed from $S$ alone coincides with the first-order wedge to machine precision. With $m=0.03$, $b=0.3$ and $e=20$, momentum traders choose $q=0.015$ against a planner's $0.0062$.

\paragraph{Interpretation.} Reversion is a public good: each dealer who leans against a displacement lowers the tail exposure of every leveraged position in the market but is paid only for his own fills. Momentum is the mirror image. Part (iv) closes the loop with the guard: the policy price of the externality is a function of the same summed response that the guard bounds, and Proposition~\ref{p2:prop:invariance} and Theorem~\ref{p2:thm:certcost} show that this quantity is identified from public data without identifying the traders. A venue that rebates reversion in proportion to $S^2/\tilde\kappa$, or a supervisor that sets pre-trade limits by $S$, prices the externality with an observable.

    Across the liquidity terciles of 30 Binance perpetuals the summed response per dollar of flow rises steeply as depth falls (Spearman correlation with log depth $-0.94$), as part (ii) predicts, but the unit-free multiple $S/|h_0|$ is the same in all three terciles (medians $46.4$, $47.3$ and $46.9$): thinner books carry more cumulative risk because their impact is larger, not because they revert more slowly. That is the pattern of part (v), which we added after seeing the comparison, so the match is a consistency check, not a test.

    \section{Evidence}\label{p2:sec:evidence}

    The empirical program tests the theory's content in order-book data. Every test fixed its specification, hashed it and committed it before its test data were opened, and every test set was opened once except where Table~\ref{p2:tab:alltests} says otherwise. Data are public: Binance USD-M futures and spot archives (one-minute klines with taker volume, depth within $\pm1$--$5\%$ of the mid at about 30-second intervals, open interest), verified against published checksums and screened by validation rules fixed in advance; Bybit trade and order-book archives; Coinbase trades; and consolidated US trades. The out-of-sample criterion is a one-sided Clark--West test \emph{and} a date-block bootstrap interval for the gain in $R^2$ that excludes zero.

    \subsection{The deepest book does not amplify}

    In BTCUSDT futures on 718 training days, the estimated one-period margin $\hat C_3$ is negative on 93.5\% of days and the median implied response peaks at impact. Local projections show no overshoot once future flow is held fixed. A registered test of whether six-hour estimates of $\hat C_3$ forecast the next hour's maximum adverse excursion beyond volatility, flow and depth controls fails ($+0.018$ points of $R^2$, interval including zero), and a volatility flag classifies tail hours far better. In spot data, a unit-free measure of transient amplification on five crash days is indistinguishable from one and from a no-feedback null, and a registered forecasting test fails. Binance's depth bands appear to show replenishment of the side being hit; on a fixed price grid built from Bybit's full order books the sign reverses, so the replenishment is an artefact of band re-centring and depletion is present, but it is not large enough to produce an overshoot. Where the theory predicts nothing, we find nothing.

    \subsection{Thin books: the summed response forecasts the size of excursions}\label{p2:sec:thin}

    Thirty Binance perpetuals were fixed from one 2024 training day before any statistic was computed, spanning BTCUSDT to symbols with under \$5M of daily volume. The amplifying regime is more frequent in thinner books (16.1\% of symbol-days with $\hat C_3>0$ in the thin tercile against 9.4\% in the thick one), although the average margin is negative everywhere. Measured feedback is larger in \emph{deeper} books: feedback and amplification are different objects.

    \begin{table}[t]\centering\small
    \begin{threeparttable}
    \caption{Registered forecasting tests in thin Binance books, 2025--2026}\label{p2:tab:thin}
    \setlength{\tabcolsep}{3pt}
    \begin{tabular}{p{0.42\textwidth}cc}
    \toprule
    Test (what changes) & $\Delta R^2$ (pts) [CI] & Clark--West $t$\\
    \midrule
    Ten thin symbols, volatility and flow controls & $+1.216$ [$+0.273$, $+7.987$] & 8.53\\
    Ten new thin symbols drawn by a seeded rule & $+0.845$ [$+0.303$, $+6.371$] & 11.66\\
    Fresh days, 20 thin symbols, price-impact controls in the baseline & $+1.583$ [$+0.588$, $+3.237$] & 9.46\\
    Fresh days, full volatility-forecasting baseline (HAR, EWMA, hour effects) & $+2.787$ [$+1.425$, $+4.944$] & 12.59\\
    \quad the guard's statistic $S_{60}$ instead of $\hat C_3$ (secondary) & $+2.098$ [$+1.165$, $+3.597$] & ---\\
    Untouched days, model frozen at its training estimates & $+4.196$ [$+1.551$, $+8.233$] & ---\\
    Untouched thin symbols & $+0.727$ [$+0.377$, $+1.428$] & ---\\
    Untouched thin symbols, beyond the baseline plus VPIN & $+0.562$ [$+0.274$, $+1.159$] & ---\\
    \bottomrule
    \end{tabular}
    \begin{tablenotes}\footnotesize\item Outcome: the next hour's maximum adverse excursion; candidate estimated from the previous six hours. Date-block bootstrap intervals (97.5\% for the fresh-day and untouched-day primaries, 95\% otherwise). Every row is a registered primary except the $S_{60}$ row. Full tables and secondaries in the full version.\end{tablenotes}
    \end{threeparttable}
    \end{table}

    The result holds in every registered forecasting test on thin Binance books (Table~\ref{p2:tab:thin}). It strengthens as the baseline gets harder: beyond a full volatility-forecasting baseline that already explains 18.6\% of the variance, $\hat C_3$ adds $2.8$ points of $R^2$, the volatility benchmarks themselves add nothing beyond the simpler controls, and the summed response $S$, the guard's own statistic, adds $2.1$ points. The decomposition matches the theory: a larger momentum term predicts larger excursions, a larger reversion term smaller ones, and the depth term nothing; most of the variation is in impact-weighted reversion, as $S=\tilde\kappa/(m-q\tilde\kappa)$ implies. It is not price impact (with four impact controls it keeps about 80\% of its gain, and impact alone forecasts nothing), not tick size (it survives tick controls, the removal of days at a non-training tick and of any single symbol; unit-free versions work on fresh days), not order-flow toxicity (it keeps about 73\% of its gain beyond VPIN, a comparison made post hoc and then confirmed in a registered test), and not order-flow persistence (a placebo with the same flow persistence but broken alignment with prices adds nothing). In economic terms, within a volatility quintile the top quintile of $\hat C_3$ has next-hour excursions about 26 basis points larger than the bottom quintile, enough to roughly double the hourly liquidation probability of a $20\times$ position.

    \subsection{Other venues}

    The result belongs to thin markets, not to one venue, but its size depends on the venue. In 23 thin US stocks a registered test passes at about a fifth of the Binance size ($+0.50$ points beyond the full baseline), carried again by the reversion term; mid-size stocks show a gain too. On Bybit the statistic is absent with trades alone (a registered failure) and, with order-book depth, meets its registered criterion with an effect about one eightieth of Binance's; on Coinbase spot the registered test fails. On both crypto venues where it fails, thin pairs trade sparsely (11\% of Bybit minutes have no trade, against 0.1\% on Binance), which leaves the series the statistic is estimated from much noisier.

    \subsection{The guard as a risk control}

    The guard's theorems concern worst-case loss under a certified operator; whether its statistic helps size positions in practice is a separate question, and the answer is: modestly, and not instead of volatility. A six-hour estimate scaled into a volatility limit does not improve it (a registered failure), and limits built on the estimate alone are worse than a static limit. With an operator certified from three days of data, the scaled limit lowers tail loss per unit of position by 2.4\% (a narrow registered pass). Used as the theorem uses it, as a one-sided throttle, the certified guard lowers tail loss against a static limit by 8\% on untouched days and 23\% on untouched symbols (registered passes), but not against a volatility cap at matched activity. Over 18-hour holding periods, an uncapped version failed its registered capacity test; a capped throttle of the kind Rule~15c3-5 requires, registered after that failure, lowers the forced-liquidation rate of a $20\times$ position by 7.4 percentage points against a volatility cap, mostly because it admits 11\% less exposure; scaled to the same mean exposure the two rules are indistinguishable.\label{p2:sec:p24}

    \subsection{What the effect is not: the mechanism tests}\label{p2:sec:p22}

    Two registered tests asked whether the forecasting content runs through forced liquidations, measured by five-minute open-interest contractions with adverse price moves. Both fail: $\hat C_3$ does not forecast them, and they mediate none of its forecasting gain ($-6.8\%$). A leverage interaction, specified in advance but not hash-registered, also fails. A follow-up registered after these failures finds that $\hat C_3$ forecasts price persistence over the next 15 to 60 minutes and that this persistence accounts for nearly all of its forecasting gain, but its registered placebo clause fails, so the mediation is not established.\label{p2:sec:p25} A zero-fee natural experiment leaves reversion unchanged and lowers impact. On October 10, 2025, reported descriptively, the summed response was elevated all day but did not warn of anything volatility missed in the final hours. We therefore read the evidence as showing that thin books are risky because their prices revert slowly relative to impact, which is what the statistic measures, not as a test of a deeper mechanism.

    \subsection{Every registered test}\label{p2:sec:alltests}
    Of twenty-five registered primaries, fifteen pass, nine fail and one was not tested because the test it depended on failed (Table~\ref{p2:tab:alltests}).
    
\begin{table}[!htbp]\centering\footnotesize\renewcommand{\arraystretch}{0.80}\setlength{\abovecaptionskip}{3pt}
\begin{threeparttable}
\caption{All pre-registered tests}\label{p2:tab:alltests}
\setlength{\tabcolsep}{3pt}
\begin{tabular}{p{0.31\textwidth}lp{0.27\textwidth}cc}
\toprule
Test (specification hash) & Sample & Primary statistic [CI] & Result & Opened\\
\midrule
BTCUSDT futures, MAE $\mid\hat C_3$ (\texttt{d9521218}) & 2025--26 & $\Delta R^2$ $+0.018$ [$-0.100$, $+0.132$] & fail & 1\\
BTCUSDT spot, MAE $\mid$ amplification (\texttt{b1fd7f8e}) & 2024 & $\Delta R^2$ $-0.023$ [$-0.056$, $-0.002$] & fail & 3\\
Thin books, 10 symbols (\texttt{5910b6b4}) & 2025--26 & $\Delta R^2$ $+1.216$ [$+0.273$, $+7.987$] & pass & 2\\
New thin symbols, 10 (\texttt{6aa8e2c0}) & 2025--26 & $\Delta R^2$ $+0.845$ [$+0.303$, $+6.371$] & pass & 1\\
Fresh days, 20 thin, impact controls (\texttt{65acda36}) & 2025--26 & $\Delta R^2$ $+1.583$ [$+0.588$, $+3.237$] & pass & 1\\
Position sizing, same specification & 2025--26 & $\Delta\log{\rm TLR}$ $-0.017$ [$-0.045$, $+0.011$] & fail & 1\\
Bybit, 14 coins, depth-free & 2025--26 & $\Delta R^2$ $-0.000$ [$-0.001$, $+0.001$] & fail & 1\\
Beyond volatility forecasting, 20 thin (\texttt{71983f1a}) & 2025--26 & $\Delta R^2$ $+2.787$ [$+1.425$, $+4.944$] & pass & 1\\
Certified guard, position sizing (\texttt{b9a56af4}) & 2025--26 & $\Delta\log{\rm TLR}$ $-0.024$ [$-0.051$, $-0.0004$] & pass & 1\\
Bybit with order-book depth, 6 coins (\texttt{ec2145fa}) & 2025--26 & $\Delta R^2$ $+0.033$ [$+0.019$, $+0.053$] & pass & 1\\
Thin US stocks, 23, depth-free (\texttt{f5800a69}) & 2025--26 & $\Delta R^2$ $+0.499$ [$+0.272$, $+0.779$] & pass & 1\\
Untouched days, frozen model (\texttt{69849c96}) & 2025--26 & $\Delta R^2$ $+4.196$ [$+1.551$, $+8.233$] & pass & 1\\
Untouched days, guard vs static (same) & 2025--26 & $\Delta\log{\rm TLR}$ $-0.084$ [$-0.173$, $-0.011$] & pass & 1\\
Untouched symbols, 5 thin (\texttt{0523e476}) & 2025--26 & $\Delta R^2$ $+0.727$ [$+0.377$, $+1.428$] & pass & 1\\
Untouched symbols, beyond VPIN (same) & 2025--26 & $\Delta R^2$ $+0.562$ [$+0.274$, $+1.159$] & pass & 1\\
Untouched symbols, guard vs static (same) & 2025--26 & $\Delta\log{\rm TLR}$ $-0.256$ [$-0.497$, $-0.021$] & pass & 1\\
New venue, Coinbase, 6 thin (\texttt{3233c52b}) & 2025--26 & $\Delta R^2$ $-0.092$ [$-0.407$, $+0.158$] & fail & 1\\
Coinbase, real minus placebo (same) & 2025--26 & $+0.013$ [$-0.331$, $+0.298$] & not tested & 1\\
OI liquidation forecast, 19 thin (\texttt{2b3794e3}) & 2025--26 & $\Delta R^2_{\rm LIQ}$ $-0.795$ [$-1.529$, $-0.179$] & fail & 1\\
OI liquidation mediation, 19 thin (same) & 2025--26 & Share $-6.8\%$ [$-10.2\%$, $-3.1\%$] & fail & 1\\
Non-resetting 18h guard capacity (\texttt{d4169377}) & 2025--26 & $\Delta\overline{N}$ (k\$) $+332.8$ [$-20.4$, $+1180.8$] & fail & 1\\
Capped 18h throttle liq.\ rate (\texttt{698f9119}) & 2025--26 & $\Delta{\rm Liq}$ $-7.41$ [$-11.45$, $-4.15$] pp & pass & 1\\
Capped 18h throttle ES$_{95}$ (same) & 2025--26 & $\Delta{\rm ES}_{95}$ (k\$) $-26.4$ [$-53.1$, $-6.3$] & pass & 1\\
Persistence forecast, 19 thin (\texttt{8ebd67f3}) & 2025--26 & $\Delta R^2_{\rm pers}$ $+3.713$ [$+1.884$, $+6.035$] & pass & 1\\
Persistence mediation + placebo (same) & 2025--26 & Share $97.6\%$ [$96.6\%$, $98.3\%$] & fail & 1\\
\bottomrule
\end{tabular}
\begin{tablenotes}\scriptsize\item $\Delta R^2$ in percentage points. ``Opened'' counts executions of the test cell; the repeated openings used unchanged code and specifications (Sections~\ref{p2:sec:thin} and~\ref{p2:sec:evidence}). Intervals are 95\%, except 97.5\% for the two fresh-day primaries and the two Part 15 primaries. \end{tablenotes}
\end{threeparttable}
\end{table}

    \paragraph{How to read the passes.} The program was sequential: each registration protected one test from specification search on its own test data, but which test to run next was chosen after the earlier results were known. The passing Binance tests share training days and are not independent replications; they vary one thing at a time (symbols, days, baseline), and the US test shares none of their data. Every feature uses only minutes before the forecast hour, and a scrambling audit leaves every feature unchanged when later data are permuted. The gains are modest on top of a strong baseline, a 15\% increase in the explained variance of the next hour's excursion at most; they are not a new predictor that replaces volatility.

    \section{Tests of the information-safety results}\label{p2:sec:istests}

    The theorems of Sections~\ref{p2:sec:infosafety}--\ref{p2:sec:disclosure} concern declared contracts, which market data cannot deliver. They are therefore tested in simulations with a known truth (Table~\ref{p2:tab:simtests}).
    
\begin{table}[!tbp]\centering\footnotesize
\begin{threeparttable}
\caption{Simulation tests of the information-safety results}\label{p2:tab:simtests}
\setlength{\tabcolsep}{3pt}
\begin{tabular}{p{0.05\textwidth}p{0.36\textwidth}p{0.5\textwidth}}
\toprule
Test & Question & Result\\
\midrule
T01 & Can markets agree on early responses yet differ in risk? & Example~\ref{p2:ex:witness} in rational arithmetic: lags 0--2 equal, summed responses $1.92$ and $5.59$\\
T02 & Does Proposition~\ref{p2:prop:twoworld} bound the capacity a safe rule gives up? & Attained exactly with identical data ($0.324$ at $\beta=0.05$); respected in 9 of 9 noisy designs by the certificate rule, whose unsafe-admission rate is at most $3.2\%$\\
T03 & Is certifying risk cheaper than identifying cohorts? & Table~\ref{p2:tab:certvsid}: slope $-6.06$ in $\Delta\eta$ for identification, $+0.23$ for certification\\
T04 & Does the certificate cover at its nominal level? & 18 designs: coverage at or above nominal in all (minimum $93.6\%$ at nominal $90\%$); median excess $3$--$15\%$ with memory bound $0.6$--$0.9$, about $230\%$ at $0.97$\\
T05 & Is the corrected gate exact? & 2{,}000 cases, horizons 1--10, every sign path enumerated: per-period worst case equals $V\sum|g|$ and total-budget worst case equals $V\max|g|$ in all; the peak gain understates the per-period worst case in 1{,}800\\
T06 & Is the closed-form guard globally optimal? & 500 economies: equal to brute force over 1{,}001 levels from the state matrices in all; safe sets always intervals; 28 infeasible states flagged\\
T07 & What does disclosure add? & Table~\ref{p2:tab:discsim}\\
T08 & Does the guard depend on the channel? & Momentum learners, fixed-rule traders, depth and slow dealer reversion calibrated to $S=5$: identical guard capacity ($0.20$); one-period margins differ ($-0.34$ to $-0.20$)\\
T09 & What breaks the certificate? & Coverage at nominal 95\%: as assumed $100\%$; delayed feedback $100\%$; contrarian cohorts $100\%$; depth channel $92.8\%$; memory longer than declared $86.3\%$; feedback 20\% stronger after certification $34.7\%$\\
T10 & Are the window endpoints proved? & Certified by verified polynomial sign conditions for every window; outside the 60-digit endpoints by at most $1.5\times10^{-8}$ relative\\
T11 & What if the tap noise is correlated? & (O1) certificate covers 74--76\% (exact value $75.1\%$); $t$-certificate 94--95\%; Hotelling 100\%; dependent experiments 82--85\%\\
T12 & Should rebates use certified bounds? & No: a 10\% overstatement of $S$ costs $2.4$--$2.8\%$ of first-best welfare, a 50\% overstatement $70$--$80\%$\\
T13 & What if admission changes the dynamics? & Safe sets $[0,0.035]\cup[0.386,1]$ (closed form) and $[0,0.053]\cup[0.178,1]$ (book, certified); no certified level unsafe in 433 sampled economies\\
T14 & Does sequential probing keep its error level? & No wrong decision in 7{,}000 runs; 3.7 to 107 probes against 130 for a fixed design\\
\bottomrule
\end{tabular}
\begin{tablenotes}\footnotesize\item All tests are deterministic given their seeds and run without data (\texttt{research\_extension/information\_safety/run\_all.py} $\to$ \texttt{results/T*.json}, about six minutes). Further checks: the joint per-period and total-budget worst case equals its linear-programming value in 300 random cases (T05b), and the balance-sheet step keeps the reserve in all 17{,}551 sampled states that satisfy the guard inequality (T12b). Kind: recomputed. \end{tablenotes}
\end{threeparttable}
\end{table}

The simulations confirm what the theorems claim and locate what they do not (Table~\ref{p2:tab:simtests}). The certificate is valid where its contract holds, including with delayed feedback and contrarian cohorts, and conservative when the declared memory is long. It fails in two ways that matter for practice: when memory is longer than declared, which is the case Theorem~\ref{p2:thm:window}(iii) warns about, and when feedback strengthens after certification, in which case a stale certificate covers the new risk only a third of the time. A guard must therefore be re-certified as conditions change; it cannot be set once. The remaining tests check the extensions: the window endpoints are proved rather than computed (T10), certificates that do not assume independent taps keep their level when the taps are correlated while the (O1) certificate does not (T11), non-monotone admission is certified by branch and bound (T13), and sequential probing keeps its error level (T14).

    \section{Policy implications}\label{p2:sec:policy}

    Rule 15c3-5 requires pre-trade controls that prevent orders exceeding credit or capital thresholds \citep{p2:sec2010}; static caps and daily value-at-risk limits do not see transient dynamics. Lemma~\ref{p2:thm:guard} names the quantity a solvency-motivated control should bound, inventory times the shock envelope times the summed absolute impulse response, and peak- or eigenvalue-based limits are unsafe or blind exactly where they matter. The evidence says where such a control matters (thin books, under about \$0.4M of depth within 1\% of the mid) and what it can do with public data: correct a volatility limit modestly, not replace it. Because the summed response is certified without identifying the traders, a supervisor or venue can set limits, or price a reversion rebate in proportion to $S^2/\tilde\kappa$, from an observable. Guards should use certified upper bounds; rebates should use point estimates.

    \section{Conclusion}\label{p2:sec:conclusion}

    The theory answers four questions. A book amplifies when depletion outweighs impact-weighted reversion; cumulative transient risk is long-run impact over reversion net of momentum; the guard on the summed response is the largest safe admission rule and needs no identification of the traders; and the information that certifies a safe decision is cheap in exactly the directions in which identification is expensive. The evidence is mixed in an informative way: the deepest book does not amplify and the statistic is useless there, while in thin books it forecasts the size of adverse excursions beyond every baseline we tried, at a size that depends on the venue, and as a practical control it adds to volatility only modestly. Message-level data from a regulated market, data on liquidity-provider identities, and a market observed near the stability boundary would settle what remains open.

    \paragraph{Data, code and AI use.} All data are public and every number is produced by archived code from them; specifications, hashes and unlock counters are in the replication package. The authors used Claude (Anthropic) to write and debug analysis code, check derivations and draft and edit text; every number comes from code the authors executed, and the authors take full responsibility for the content.
    
\begin{thebibliography}{99}\setlength{\itemsep}{2pt}\small
\bibitem[21Shares(2025)]{p2:shares2025} 21Shares (2025). Record crypto liquidations amid tariff shock. Research note, October 13, 2025. \url{https://www.21shares.com/research/record-crypto-liquidations-amid-tariff-shock}.
\bibitem[Achieser(1956)]{p2:achieser1956} Achieser, N. I. (1956). \emph{Theory of Approximation}. New York: Frederick Ungar.
\bibitem[Andersen and Bondarenko(2014)]{p2:andersen2014} Andersen, T. G., and O. Bondarenko (2014). VPIN and the flash crash. \emph{Journal of Financial Markets} 17, 1--46.
\bibitem[Bouchaud et~al.(2004)]{p2:bouchaud2004} Bouchaud, J.-P., Y. Gefen, M. Potters, and M. Wyart (2004). Fluctuations and response in financial markets: The subtle nature of ``random'' price changes. \emph{Quantitative Finance} 4, 176--190.
\bibitem[Bouchaud et~al.(2018)]{p2:bouchaud2018} Bouchaud, J.-P., J. Bonart, J. Donier, and M. Gould (2018). \emph{Trades, Quotes and Prices: Financial Markets Under the Microscope}. Cambridge University Press.
\bibitem[Brunnermeier and Pedersen(2009)]{p2:brunnermeier2009} Brunnermeier, M. K., and L. H. Pedersen (2009). Market liquidity and funding liquidity. \emph{Review of Financial Studies} 22, 2201--2238.
\bibitem[Clark and West(2007)]{p2:clark2007} Clark, T. E., and K. D. West (2007). Approximately normal tests for equal predictive accuracy in nested models. \emph{Journal of Econometrics} 138, 291--311.
\bibitem[Corsi(2009)]{p2:corsi2009} Corsi, F. (2009). A simple approximate long-memory model of realized volatility. \emph{Journal of Financial Econometrics} 7, 174--196.
\bibitem[Dahleh and Diaz-Bobillo(1995)]{p2:dahleh1995} Dahleh, M. A., and I. J. Diaz-Bobillo (1995). \emph{Control of Uncertain Systems: A Linear Programming Approach}. Prentice Hall.
\bibitem[Dahleh and Pearson(1987)]{p2:dahleh1987} Dahleh, M. A., and J. B. Pearson (1987). $\ell^1$-optimal feedback controllers for MIMO discrete-time systems. \emph{IEEE Transactions on Automatic Control} 32, 314--322.
\bibitem[Donoho(1994)]{p2:donoho1994} Donoho, D.~L. (1994). Statistical estimation and optimal recovery. \emph{Annals of Statistics} 22, 238--270.
\bibitem[Easley et~al.(2012)]{p2:easley2012} Easley, D., M. M. L\'opez de Prado, and M. O'Hara (2012). Flow toxicity and liquidity in a high-frequency world. \emph{Review of Financial Studies} 25, 1457--1493.
\bibitem[Hansen and Walster(2004)]{p2:hansen2004} Hansen, E., and G. W. Walster (2004). \emph{Global Optimization Using Interval Analysis}, 2nd ed. New York: Marcel Dekker.
\bibitem[Harvey(2017)]{p2:harvey2017} Harvey, C. R. (2017). Presidential address: The scientific outlook in financial economics. \emph{Journal of Finance} 72, 1399--1440.
\bibitem[Hotelling(1931)]{p2:hotelling1931} Hotelling, H. (1931). The generalization of Student's ratio. \emph{Annals of Mathematical Statistics} 2, 360--378.
\bibitem[J.P. Morgan/Reuters(1996)]{p2:riskmetrics1996} J.P. Morgan/Reuters (1996). \emph{RiskMetrics---Technical Document}, 4th ed. New York: J.P. Morgan.
\bibitem[Karlin and Studden(1966)]{p2:karlin1966} Karlin, S., and W. J. Studden (1966). \emph{Tchebycheff Systems: With Applications in Analysis and Statistics}. New York: Interscience.
\bibitem[Kirilenko et~al.(2017)]{p2:kirilenko2017} Kirilenko, A., A. S. Kyle, M. Samadi, and T. Tuzun (2017). The Flash Crash: High-frequency trading in an electronic market. \emph{Journal of Finance} 72, 967--998.
\bibitem[Krein and Nudel'man(1977)]{p2:krein1977} Krein, M. G., and A. A. Nudel'man (1977). \emph{The Markov Moment Problem and Extremal Problems}. Translations of Mathematical Monographs 50. Providence, RI: American Mathematical Society.
\bibitem[Le Cam(1986)]{p2:lecam1986} Le Cam, L. (1986). \emph{Asymptotic Methods in Statistical Decision Theory}. New York: Springer.
\bibitem[Milanese and Vicino(1991)]{p2:milanese1991} Milanese, M., and A. Vicino (1991). Optimal estimation theory for dynamic systems with set membership uncertainty: An overview. \emph{Automatica} 27, 997--1009.
\bibitem[Nosek et~al.(2018)]{p2:nosek2018} Nosek, B. A., C. R. Ebersole, A. C. DeHaven, and D. T. Mellor (2018). The preregistration revolution. \emph{Proceedings of the National Academy of Sciences} 115, 2600--2606.
\bibitem[Obizhaeva and Wang(2013)]{p2:obizhaeva2013} Obizhaeva, A. A., and J. Wang (2013). Optimal trading strategy and supply/demand dynamics. \emph{Journal of Financial Markets} 16, 1--32.
\bibitem[Rump(2010)]{p2:rump2010} Rump, S. M. (2010). Verification methods: Rigorous results using floating-point arithmetic. \emph{Acta Numerica} 19, 287--449.
\bibitem[SEC(2010)]{p2:sec2010} U.S. Securities and Exchange Commission (2010). Risk management controls for brokers or dealers with market access. Release No.\ 34-63241.
\bibitem[Shleifer and Vishny(1992)]{p2:shleifer1992} Shleifer, A., and R. W. Vishny (1992). Liquidation values and debt capacity: A market equilibrium approach. \emph{Journal of Finance} 47, 1343--1366.
\bibitem[Sornette et~al.(2023)]{p2:sornette2023} Sornette, D., S. Lera, J. Lin, and K. Wu (2023). Non-normal interactions create socio-economic bubbles. \emph{Communications Physics} 6, 261.
\bibitem[Tu et~al.(2018)]{p2:tu2018} Tu, S., R. Boczar, and B. Recht (2018). Minimax lower bounds for $\mathcal H_\infty$-norm estimation. Working paper, arXiv:1809.10855.
\bibitem[Vidyasagar(1986)]{p2:vidyasagar1986} Vidyasagar, M. (1986). Optimal rejection of persistent bounded disturbances. \emph{IEEE Transactions on Automatic Control} 31, 527--534.
\bibitem[van Waarde et~al.(2020)]{p2:vanwaarde2020} van Waarde, H. J., J. Eising, H. L. Trentelman, and M. K. Camlibel (2020). Data informativity: A new perspective on data-driven analysis and control. \emph{IEEE Transactions on Automatic Control} 65, 4753--4768.
\end{thebibliography}

    \section{Appendix: Proofs}\label{p2:app:proofs}
    
\paragraph{Proposition~\ref{p2:prop:jury}.} With one cohort, $J=\begin{psmallmatrix}a&\kappa q\\ \eta a&1-\eta+\eta\kappa q\end{psmallmatrix}$, so ${\rm tr}J=a+1-\eta+\eta\kappa q$ and $\det J=a(1-\eta)$. The Jury conditions for a real $2\times2$ matrix are $|\det J|<1$ and $1\pm{\rm tr}J+\det J>0$. Substituting $a=1-\kappa\varphi_f$: $1-{\rm tr}J+\det J=\eta(1-a-\kappa q)=\eta\kappa(\varphi_f-q)$, positive iff $q<\varphi_f$; and $1+{\rm tr}J+\det J=(4-2\eta)-\kappa[(2-\eta)\varphi_f-\eta q]$, positive iff the stated bound on $\kappa$ holds (automatically when $(2-\eta)\varphi_f\le\eta q$).

\paragraph{Proposition~\ref{p2:prop:invariance}.} Equal merged measures give equal $G$, hence equal $H$, and the impulse coefficients are the Laurent coefficients of $H$ (Chapter~\ref{ch:p1}, Theorem on merged equivalence). Under $z\mapsto Tz$, $c^\top J^kb$ is unchanged because $CT^{-1}(TJT^{-1})^kTb=CJ^kb$, while the row norm $\|CT^{-1}(TJT^{-1})^k\|$ depends on $T$. The displayed price path is the convolution of inputs with $\{g_k\}$, and its worst case over the $\ell_\infty$ ball is $\bar u\sum|g|$ by H\"older.

\paragraph{Proposition~\ref{p2:prop:inventory_S}.} (i) Taking $z$-transforms of $y_{i,t}=(1-\eta_i)y_{i,t-1}+\eta_i x_t$, $I_{i,t}=(1-\delta_{\rm inv})I_{i,t-1}+u_{i,t}$ and $u_{i,t}=q_i y_{i,t}-\psi_i I_{i,t-1}$ gives $Y_i(z)=\frac{\eta_i z}{z-1+\eta_i}X(z)$ and $(z-1+\delta_{\rm inv}+\psi_i)I_i(z)=q_i z Y_i(z)$, whence $U_i(z)=\frac{z-1+\delta_{\rm inv}}{z-1+\delta_{\rm inv}+\psi_i}q_i Y_i(z)$ and $G_{\rm inv}(z)=\sum_i U_i(z)/X(z)$ has the stated form. Evaluating at $z=1$ when the augmented operator is Schur stable gives $G_{\rm inv}(1)=\sum_i \frac{\delta_{\rm inv}}{\delta_{\rm inv}+\psi_i}q_i=\Phi_{\rm eff}(\psi)$ and $S(\psi)=H_{\rm inv}(1)=\kappa/(1-a-\kappa\Phi_{\rm eff}(\psi))$. (ii) Write $r_i=1-\eta_i$ and $\rho_i=1-\delta_{\rm inv}-\psi_i$. When $0\le\psi_i\le 1-\delta_{\rm inv}$ and $\delta_{\rm inv}\ge\eta_i$, we have $0\le\rho_i\le 1-\delta_{\rm inv}\le r_i<1$. For $r_i>\rho_i$, the $m$-th open-loop tap is
$k_{i,m}=q_i\eta_i[r_i^m-\psi_i(r_i^m-\rho_i^m)/(r_i-\rho_i)]
=q_i\eta_i[(1-v_i)r_i^m+v_i\rho_i^m]$, where
$v_i=\psi_i/(r_i-\rho_i)=\psi_i/(\delta_{\rm inv}-\eta_i+\psi_i)\in[0,1]$.
It is nonnegative. If $r_i=\rho_i$, the assumptions force $\psi_i=0$ and the tap is $q_i\eta_i r_i^m\ge0$ (with the usual $r_i^0=1$ convention).
Nonnegative open-loop taps $k_m(\psi)\ge0$ with $a+\kappa\sum_m k_m(\psi)=a+\kappa\Phi_{\rm eff}(\psi)<1$ generate a nonnegative closed-loop renewal sequence $h_t(\psi)\ge0$, so $\|h(\psi)\|_1=\sum_{t\ge0}h_t(\psi)=S(\psi)$.

\paragraph{Theorem~\ref{p2:thm:radius}.} (i) $\sup_t\|CJ^tz_0\|\le\sup_t\|CJ^t\|\|z_0\|$, with equality approached by aligning $z_0$ with the top right singular vector of $CJ^{t^*}$ at the maximizing $t^*$. (ii) The Kreiss matrix theorem with Spijker's sharp constant gives $\mathcal K(J)\le\sup_t\|J^t\|\le en\mathcal K(J)$; substitute in (i). (iii) Take $J_\theta=\begin{psmallmatrix}\rho&\theta\\0&\rho/2\end{psmallmatrix}$; then $\rho(J_\theta)=\rho$ and $(1,0)J_\theta e_2=\theta$, so $M_C\ge|\theta|$ and $r^*_{\rm exact}\le B/|\theta|\to0$.

\paragraph{Proposition~\ref{p2:prop:driftguard}.} (i) Positivity gives $|c^\top J^tz_0|\le c^\top J^t\bar z_0$, with equality for memories aligned with $x_0$ (either sign if $x_0=0$). Summing the original memory recursion, without rescaling its initial condition, gives $\sum_{t\ge0}y_{i,t}=y_{i,0}/\eta_i+\sum_{t\ge1}x_t$. Hence $(1-a-\alpha\kappa\Phi)\sum_{t\ge1}x_t=ax_0+\alpha\kappa\sum_iq_iy_{i,0}/\eta_i$, proving \eqref{p2:eq:driftenv}. (ii) Superposition and positivity give the expression in \eqref{p2:eq:driftpeak}; shocks and memories with the common sign of $x_0$ attain it. Taking the maximum over the same time index proves sharpness for fixed parameters. (iii) If $|x_t|,|y_{i,t}|\le L$, then $|x_{t+1}|\le aL+\alpha\kappa\Phi^+L+\kappa\bar u\le L$. Each subsequent memory is a convex combination of two values bounded by $L$. Induction proves the invariant box. Rearranging gives \eqref{p2:eq:alphabox}. For $\bar u>0$, its endpoint automatically has $a+\alpha\kappa\Phi^+<1$. No general maximality or positive capacity cost is inferred by adding separately maximized cumulative and pointwise quantities.

\paragraph{Theorem~\ref{p2:thm:window}.} The recursion $h_{t+1}=ah_t+\kappa\sum_{m\le t}k_mh_{t-m}$ with $h_0=\kappa\neq0$ gives $k_t$ from $h_0,\ldots,h_{t+1}$, so the window determines $k_0,\ldots,k_{n-1}$ and nothing more about $M$: any nonnegative measure on $[0,\bar r]$ with these moments, written as atoms $(r_i,w_i)$ or as a limit of such, is the merged measure of a positive economy with $\eta_i=1-r_i$ and $q_i=w_i/\eta_i$ that produces the same window. (i) The compatible measures form a convex, weak-* compact set, and $\Phi=\int f\,dM$ is linear and continuous on it ($f$ is continuous on $[0,\bar r]$), so its range is a closed interval $[\Phi_-,\Phi_+]$. Since $f^{(n)}(r)=n!/(1-r)^{n+1}>0$ on $[0,\bar r]$, $(1,r,\ldots,r^{n-1},f)$ is an extended complete Chebyshev system, and by the Markov--Krein theorem \citep[Theorems IV.1.1 and IV.6.1]{p2:krein1977} the maximum is attained at the upper and the minimum at the lower principal representation, which have at most $\lfloor n/2\rfloor+1$ atoms. The theorem applies to moment vectors in the interior of the moment space, where both principal representations exist and are unique; a vector on the boundary has exactly one representing measure, so the interval collapses to a point. $S=\kappa/(1-a-\kappa\Phi)$ is increasing in $\Phi$ on $\Phi<(1-a)/\kappa$, which gives the interval of summed responses, unbounded when $\Phi_+\ge(1-a)/\kappa$ because by convexity every $\Phi\in[\Phi_-,\Phi_+]$ is attained. (ii) If $\mu,\nu$ share moments up to order $n-1$, then $\int f\,d(\mu-\nu)=\int(f-p)\,d(\mu-\nu)$ for every polynomial $p$ of degree $n-1$, so $|\int f\,d(\mu-\nu)|\le2k_0\inf_p\|f-p\|_\infty=2k_0E_{n-1}$. With $r=\bar r(1+x)/2$, $f=(2/\bar r)/(x_0-x)$ on $[-1,1]$ with $x_0=2/\bar r-1>1$, and Chebyshev's formula gives $E_m(1/(x_0-x))=\varrho^m/(x_0^2-1)$ \citep{p2:achieser1956}; we checked it against a linear-programming minimax fit for $m\le8$. Expanding $x_0=1+2(1-\bar r)/\bar r$ gives the asymptotics. (iii) Let the moment vector $c$ be interior to the moment cone of $[0,1]$. Interior points have an open neighbourhood in the cone, and $m(r_0)=(1,r_0,\ldots,r_0^{n-1})$ has norm at most $\sqrt n$, so for some $\epsilon>0$, $c-\epsilon m(r_0)$ is interior for every $r_0\in[0,1]$ and is represented by its lower principal representation, which has no atom at $1$. Adding $\epsilon\delta_{r_0}$ gives a measure on $[0,1)$ with moments $c$ and $\Phi\ge\epsilon/(1-r_0)\to\infty$ as $r_0\to1$; by convexity every larger $\Phi$ below $(1-a)/\kappa$ is attained by a stable economy, so $S$ is unbounded.

\paragraph{Theorem~\ref{p2:thm:certcost}.} (i) The flow impulse enters the stacked state as $b=\kappa(1,\eta_1,\ldots,\eta_K)^\top$, and $h_t=e_x^\top J^tb\le\|J^tb\|_\infty\le\|J\|_\infty^t\|b\|_\infty$ because $J\ge0$. The price row of $J$ sums to $a+\kappa\Phi\le\bar s$ and cohort row $i$ to $1-\eta_i(1-a-\kappa\Phi)\le1-(1-\bar r)(1-\bar s)=\bar\lambda$, so $h_t\le\kappa\bar\lambda^t$ and $\sum_{t\ge T}h_t\le\kappa\bar\lambda^T/(1-\bar\lambda)$. $\sum_{t<T}\bar h_t$ is Gaussian with mean $\sum_{t<T}h_t$ and variance $T\sigma^2/N$, so $\sum_{t<T}\bar h_t+z_{1-\gamma}\sigma\sqrt{T/N}\ge\sum_{t<T}h_t$ with probability $1-\gamma$, and adding the tail bound gives validity. On the two-sided event of probability $1-2\gamma$, $\hat S^+-S\le2z_{1-\gamma}\sigma\sqrt{T/N}+\kappa\bar\lambda^T/(1-\bar\lambda)$; $T_\epsilon$ makes the second term at most $\epsilon/2$ and the stated $N$ makes the first at most $\epsilon/2$. (ii) Let $\theta_1$ raise the cohort weight of $\theta_0$ by $\Delta$ chosen so that $S_1-S_0=2\epsilon$; then $\|h_1-h_0\|_2=2C_0\epsilon(1+o(1))$. Validity at $\theta_1$ gives $P_1(\hat S\ge S_1)\ge1-\gamma$, hence $P_0(\hat S-S_0\ge2\epsilon)\ge1-\gamma-{\rm TV}(P_0,P_1)$. For $N$ independent Gaussian responses, ${\rm KL}=N\|h_1-h_0\|_2^2/(2\sigma^2)$, and Pinsker's inequality gives ${\rm TV}\le\sqrt N\,C_0\epsilon(1+o(1))/\sigma$. An overshoot probability below $1/2$ therefore needs ${\rm TV}>1/2-\gamma$, which is the stated bound. (iv) Write $M=\sum_iq_i(1-r_i)\delta_{r_i}$ for the merged measure, $k_m=\int r^m dM$ for the open-loop kernel and $\mathcal M_0$ for the measures on $[0,\bar r]$ of mass at most $B_0=(1-\bar r)(\bar s-a)/\kappa$; economies with merged measure in $\mathcal M_0$ lie in the class, and finite-atom measures are weak-$*$ dense in $\mathcal M_0$. \emph{Step 1 (closed loop to open loop).} The closed-loop transforms satisfy $H=\kappa/(1-az-\kappa zK)$, so $H_1-H_0=\kappa^{-1}zH_1H_0(K_1-K_0)$ and, by Young's inequality, $\|h_1-h_0\|_2\le\kappa^{-1}S_1S_0\|k_1-k_0\|_2\le C_1\|k_1-k_0\|_2$ on the class, for any number of recorded responses. \emph{Step 2 (two points).} If two economies have $S_1>S_0+\epsilon$, validity at the first and the overshoot bound at the second make $\{\hat S^+\le S_0+\epsilon\}$ have probability at least $1-\gamma$ under $P_0$ and at most $\gamma$ under $P_1$, so the total variation is at least $1-2\gamma$; for Gaussian responses this needs $\|h_1-h_0\|_2\ge2z_{1-\gamma}\sigma/\sqrt N$. Hence $\omega_S(2z_{1-\gamma}\sigma/(C_1\sqrt N))\le\epsilon$, where $\omega_S(\delta)$ is the largest difference in $S$ between two economies with $\|k_1-k_0\|_2\le\delta$. Since $\partial S/\partial\Phi=\kappa^2/(1-a-\kappa\Phi)^2\ge\kappa^2$, $\omega_S\ge\kappa^2\omega_\Phi$. \emph{Step 3 (duality).} For a depth $L$, $\omega_\Phi(\delta)$ over $\mathcal M_0$ with the first $L$ kernel values equals $\min_{c\in\mathbb R^L}\{\delta\|c\|_2+B_0\psi(g-p_c)\}$, with $g(r)=1/(1-r)$, $p_c(r)=\sum_{t<L}c_tr^t$ and $\psi(h)=(\sup h)_++(-\inf h)_+$ on $[0,\bar r]$: the inequality $\le$ is direct, and equality follows from Sion's minimax theorem because $\mathcal M_0^2$ is weak-$*$ compact \citep{p2:donoho1994}; the modulus for all kernel values is the limit as $L\to\infty$, by compactness. \emph{Step 4 (a pole forces large coefficients).} If $\|c\|_2\le R$ and $\eta=\sup_{[0,\bar r]}|g-p_c|$, then $\log(1/\eta)\le K_0\max\{R^2,(1-\bar r)^{-1}\}$ with $K_0=8\log6\sqrt{1-\bar r}\log\rho_*$, $\rho_*=(1+\sqrt{1-\bar r})^2/\bar r$. Indeed $e=g-p_c$ is analytic on the disc of radius $\rho=1-s/2$, $s=\min\{1/(4R^2),(1-\bar r)/4\}$, with $|e(1-s)|\ge1/(2s)$ and $|e|\le3/s$ on $|z|=\rho$ because $|p_c(z)|\le R(1-|z|)^{-1/2}$; the two-constants theorem on the disc slit along $[0,\bar r]$, with the harmonic measure of the slit bounded below by an annulus comparison through the Joukowski map, gives the claim. \emph{Step 5.} For every $c$, either $\eta\ge\delta\sqrt{\log(1/\delta)}$, or $\log(1/\eta)>\tfrac12\log(1/\delta)$ and Step 4 forces $\|c\|_2\ge\sqrt{\log(1/\delta)/(2K_0)}$ once $\log(1/\delta)\ge2K_0/(1-\bar r)$; so $\omega_\Phi(\delta)\ge\min\{B_0,(2K_0)^{-1/2}\}\delta\sqrt{\log(1/\delta)}$ for small $\delta$, and Step 2 gives $N\ge c\,\sigma^2\epsilon^{-2}\log(1/\epsilon)$. (v) The identification rates are Chapter~\ref{ch:p1}'s Theorems on the $K$-cohort lower bound, the matching upper bound and general memory.

\paragraph{Lemma~\ref{p2:lem:sound}.} On the event that $\mathcal C(D)$ contains the true economy, the supremum over $\mathcal C(D)$ dominates the true worst-case loss for every action, including one selected from $D$; the event has probability at least $1-\gamma$.

\paragraph{Theorem~\ref{p2:thm:robustguard}.} (i) For $c_I\ge0$, $\alpha\mapsto S(\alpha)=\kappa/(1-a-\kappa\alpha\Phi)$ is increasing, and so is the worst-case loss $|I|\bar uS(\alpha)$, which is also increasing in $\Phi$; the robust constraint is therefore the constraint at $\Phi^+$, and $|I|\bar u\kappa/(1-a-\kappa\alpha\Phi^+)\le\mathcal B$ is equivalent to $\alpha\le c_I/(\kappa\Phi^+)$. (ii) is Lemma~\ref{p2:lem:sound}. (iii) At $\Phi=\Phi^+$ the constraint binds at $\alpha^*$ when $\alpha^*<1$. (iv) At $\alpha=0$, $S=\kappa/(1-a)$. (v) Upper: $\alpha^*$ is Lipschitz in $\Phi^+$ with constant $c_I/(\kappa\Phi_{\min}^2)$ and $\Phi^+$ is Lipschitz in $\hat S^+$ with constant $1/S_{\min}^2$, so Theorem~\ref{p2:thm:certcost}(i) with $T\asymp\log N$ gives the rate. Lower: take two single-cohort economies in the class whose weights differ by $\Delta=c\,\sigma/\sqrt N$; their oracle capacities differ by a constant times $\Delta$, and ${\rm TV}\le\sqrt N\|\partial h/\partial q\|_2\Delta/(2\sigma)$, which is below $(1-\beta)/2$ for small $c$; Proposition~\ref{p2:prop:twoworld} gives an expected loss of at least $(1-\beta)/2$ times the capacity difference.

\paragraph{Theorem~\ref{p2:thm:disclosure}.} (i) If $O'$ reveals everything $O$ reveals, every economy compatible with $O'$ data is compatible with $O$ data, so suprema over $\Theta_{O'}$ are at most those over $\Theta_O$, and Theorem~\ref{p2:thm:robustguard}(i) is monotone in $\Phi^+$. (ii) $\alpha^*$ is strictly decreasing in $\Phi^+$ when $c_I>0$ and $\alpha^*<1$, and $\Phi^+_{O'}<\Phi^+_O$ exactly when no economy with $\Phi=\Phi^+_O$ remains compatible (the supremum is attained, by Theorem~\ref{p2:thm:window}(i)). (iii) follows from Theorem~\ref{p2:thm:window}(iii). (iv) If $O$ identifies $\Phi$, $\Phi^+_O=\Phi$ and no contract can raise $\alpha^*$. (v) The delta method applied to $(w,r)\mapsto w/(1-r)$, whose gradient is $(1-r)^{-1}(1,\Phi_i)$, gives the standard error. If each certificate is valid at level $1-\gamma/2$, their minimum is valid at level $1-\gamma$ by the union bound and is no larger than either, so the admitted level is at least the larger of the two admitted levels.

